\chapter{演算法分析與漸近符號}
\label{ch:analysis}

% 本章的圖共用的樣式
\tikzset{
    sortedcell/.style={cell, fill=coreFill, draw=coreStroke},
    growth/.style={thick, smooth},
    glabel/.style={font=\scriptsize, anchor=west},
}
% 插入排序追蹤圖的一列：\insrow{y}{左邊的說明}{數值串列}{已排好的格數}{剛插入的格子}
\newcommand{\insrow}[5]{%
    \node[anchor=east, font=\small] at (0.3, #1) {#2};
    \foreach \v [count=\k] in {#3} {
        \ifnum\k=#5\relax
            \node[hlcell] at (\k*0.9, #1) {\v};
        \else\ifnum\k>#4\relax
            \node[cell] at (\k*0.9, #1) {\v};
        \else
            \node[sortedcell] at (\k*0.9, #1) {\v};
        \fi\fi
    }
}

系辦公室收到一批暑期實習的申請，每份申請寫著一個九位數的學號。 有些同學怕沒
送到，同一份申請送了兩次，助理想先知道這批申請裡有沒有重複的學號。 第一種做
法是拿每一份去和其他每一份比對學號；第二種做法是準備一張很大的表，學號每出現一
次就在表上對應的格子打勾，打勾前發現格子已經勾過，就是重複。 申請從 $1000$ 份變
成 $8000$ 份時，兩種做法各要多做幾倍的工作？ 往下讀之前，先估計一下。

這一章講怎麼回答這種問題。 我們不用碼錶量秒數，改成數一個演算法做了多少個基本
步驟，把步數寫成輸入大小 $n$ 的函數，再用\textbf{漸近符號}（asymptotic notation）
$O$、$\Omega$、$\Theta$ 只留下 $n$ 很大時決定快慢的那一項。 中間會用插入排序當例
子，示範三件之後每一章都會用到的事：逐行算出執行次數、用迴圈不變量說明程式為什麼
對、分別討論最好、最壞和平均的情況。

這一章用到《C 程式設計基礎》第~9~章的指標參數（用指標把計數帶回呼叫的地方）和
第~19~章的 \texttt{malloc}、\texttt{calloc}、\texttt{free}，用到時會先在框裡複習。
數學上只需要高中的指數、對數和數列求和；本書用到的求和公式整理在附錄~A。

\section{演算法和它的成本}
\label{sec:ana-cost}

\textbf{演算法}（algorithm）是一串定義清楚的步驟，把\textbf{輸入}（input）變成
\textbf{輸出}（output）。 「定義清楚」的意思是每一步都沒有歧義，任何人照著做都
得到一樣的結果；而且對每一個合法的輸入，它都會在有限步之後停下來，給出正確的輸
出。 一個\textbf{問題}（problem）規定了輸入和輸出的關係，例如「給 $n$ 個學號，回
答有沒有兩個相同」；把具體的一批學號交給它，叫做這個問題的一個\textbf{實例}
（instance）。

\subsection{數步驟，不量秒數}

同一支程式在不同電腦、不同編譯器上跑，秒數可以差好幾倍；同一台電腦上，秒數也會
受到其他程式影響。 我們想要的是一個只和演算法本身有關的量，所以改成數步驟。 數
法建立在一個簡化的電腦模型上，叫\textbf{隨機存取機}（random-access machine，簡稱
RAM 模型）：

\begin{itemize}
\item 一個\textbf{基本運算}花固定的時間：整數的加減乘除和取餘數、比較兩個數、
指定一個變數、讀寫陣列的一格、呼叫函式和返回。
\item 讀寫記憶體的任何一個位置都花一樣的時間，不管位置在哪裡。
\item 一次只做一個運算，沒有平行。
\item 一個整數放得進一個「字組」（word），大小足以存下 $n$ 這個量級的數。 不能把
任意大的數塞進一個字組，假裝運算還是固定時間；第~\ref{sec:ana-input-size}~節會
看到這件事為什麼重要。
\end{itemize}

這個模型和真正的電腦不完全一樣（真正的電腦有快取，讀相鄰的記憶體比較快），但用它
算出來的步數，在 $n$ 變大時的變化趨勢通常和實際的秒數一致。 這本書說「時間」
或「執行時間」（running time），指的都是在這個模型下的步數。

步數會隨輸入變多而變多，所以要先說好怎麼量輸入的大小。 對大部分問題，\textbf{輸
入大小}（input size）就是資料的個數：排序 $n$ 個數，大小是 $n$；圖的問題要用兩個
數，點數 $n$ 和邊數 $m$（第~21~章）。 有些問題的輸入只有一兩個整數，這時要改用
位元數來量，第~\ref{sec:ana-input-size}~節再講。

\subsection{例子：找重複的學號}

章首的兩種做法寫成虛擬碼如下。 學號都在 $[\mathit{base}, \mathit{base} + M)$ 這
個範圍裡，第二種做法用一個長度 $M$ 的陣列 $\mathit{seen}$ 打勾。

\begin{pseudo}[label=alg:dup-pairs]{兩兩比較}
\Function{Has-Dup-Pairs}{$A, n$}
    \For{$i \gets 1$ \To $n - 1$}
        \For{$j \gets i + 1$ \To $n$}
            \If{$A[i] = A[j]$}
                \State \Return \TRUE
            \EndIf
        \EndFor
    \EndFor
    \State \Return \FALSE
\EndFunction
\end{pseudo}

\begin{pseudo}[label=alg:dup-marks]{在表上打勾}
\Function{Has-Dup-Marks}{$A, n, \mathit{base}, M$}
    \State 建立陣列 $\mathit{seen}[0..M-1]$，每一格都設成 \FALSE
    \For{$i \gets 1$ \To $n$}
        \State $k \gets A[i] - \mathit{base}$
        \If{$\mathit{seen}[k]$}
            \State \Return \TRUE
        \EndIf
        \State $\mathit{seen}[k] \gets \TRUE$
    \EndFor
    \State \Return \FALSE
\EndFunction
\end{pseudo}

沒有重複的時候，第一種做法把每一對都比過一次。 $i = 1$ 時比 $n - 1$ 次，$i = 2$
時比 $n - 2$ 次……最後 $i = n - 1$ 時比 $1$ 次，總共
\[
(n - 1) + (n - 2) + \cdots + 1 = \frac{n(n-1)}{2}
\]
次。 第二種做法每個學號查一次表，共 $n$ 次，但開始前要把 $M$ 格都設成
\FALSE，這也要 $M$ 步。

下面的 C 程式把兩種做法各寫成一個函式，用指標參數把「比較了幾次」「查了幾次表」
帶回 \texttt{main}。 學號用 \texttt{BASE + i * 7919 \% RANGE} 產生，$7919$ 和
$10^6$ 沒有公因數，所以 $n \le 10^6$ 時產生的學號都不同，這是兩種做法最花時間的
情況（找到重複就可以提早結束）。

\begin{cbookrecap}[第~9~章 9.2 節、第~19~章 19.3 節]
函式只能回傳一個值。 想多帶回一個數，就多傳一個指標參數，在函式裡透過
\texttt{*cmp} 改它指到的變數，例如下面的 \texttt{has\_dup\_pairs(id, n, \&a)}。

\texttt{malloc(n * sizeof *id)} 向系統要一塊能放 $n$ 個 \texttt{int} 的記憶體，
要不到時回傳 \texttt{NULL}。 \texttt{calloc(count, size)} 一樣要記憶體，但會把整
塊設成 $0$。 用完要 \texttt{free}，同一塊不能 \texttt{free} 兩次。 本書所有的程式
都用 \texttt{-fsanitize=address} 檢查過，忘了 \texttt{free} 或是讀寫超出範圍都會
被抓出來。
\end{cbookrecap}

\begin{example}[label=ex:dup-check]{c}{(dup\_check.c) 找重複學號的兩種做法}
#include <stdio.h>
#include <stdlib.h>

#define BASE 115000000          // 學號的最小值
#define RANGE 1000000           // 學號落在 [BASE, BASE + RANGE)

// 方法一：兩兩比較。*cmp 記錄比較了幾次。
int has_dup_pairs(const int *id, int n, long long *cmp)
{
    *cmp = 0;
    for (int i = 0; i < n; i++) {
        for (int j = i + 1; j < n; j++) {
            (*cmp)++;
            if (id[i] == id[j]) {
                return 1;
            }
        }
    }
    return 0;
}

// 方法二：用一個長度 RANGE 的陣列做記號。*look 記錄查了幾次記號。
int has_dup_marks(const int *id, int n, long long *look)
{
    char *seen = calloc(RANGE, 1);   // 全部設成 0，代表都還沒看過
    if (seen == NULL) {
        return -1;
    }
    int found = 0;
    *look = 0;
    for (int i = 0; i < n && !found; i++) {
        (*look)++;
        int k = id[i] - BASE;
        if (seen[k]) {
            found = 1;
        }
        seen[k] = 1;
    }
    free(seen);
    return found;
}

int main(void)
{
    printf("n pairs marks result\n");
    for (int n = 1000; n <= 8000; n *= 2) {
        int *id = malloc(n * sizeof *id);
        if (id == NULL) {
            return 1;
        }
        for (int i = 0; i < n; i++) {   // 產生 n 個互不相同的學號
            id[i] = BASE + (int)((i * 7919LL) % RANGE);
        }
        long long a, b;
        int r1 = has_dup_pairs(id, n, &a);
        int r2 = has_dup_marks(id, n, &b);
        printf("%d %lld %lld %s\n", n, a, b,
               (r1 == 0 && r2 == 0) ? "no-dup" : "dup");
        free(id);
    }
    return 0;
}
\end{example}
\iosample{}{n pairs marks result \\ 1000 499500 1000 no-dup \\ 2000 1999000 2000 no-dup \\ 4000 7998000 4000 no-dup \\ 8000 31996000 8000 no-dup}

$n$ 每變成兩倍，兩兩比較的次數大約變成四倍，查表的次數剛好變成兩倍。 從 $1000$ 份
到 $8000$ 份，第一種做法多做了大約 $64$ 倍的工作，第二種只多了 $8$ 倍。 章首的問
題答案就是這兩個數。

第二種做法的代價在別的地方：不管 $n$ 多小，它都要先準備 $M = 10^6$ 格的表，花
$M$ 個位元組的記憶體，\texttt{calloc} 把它們設成 $0$ 也要時間，所以它的總步數比
較像 $n + M$。 學號如果是任意的 $64$ 位元整數，$M$ 會大到根本開不出這張表。 這時
要改用第~12~章的雜湊，或是先排序（第~14~章）再比相鄰的兩個。 一個演算法好不
好，要同時看時間、記憶體，以及它對輸入做了什麼假設。

\section{本書的虛擬碼和 C 程式}
\label{sec:ana-conventions}

這本書每個演算法都會給兩種寫法：虛擬碼和 C 程式。 虛擬碼的寫法和多數演算法教科
書、考卷上的寫法相近，用來講道理、算次數、手算追蹤；C 程式用來實際執行，驗證書上
說的結果。 兩種寫法有幾個地方不同，下表整理在一起，之後各章不再重複。 表裡只對照寫法，$i$ 的值兩邊相同；陣列編號差 $1$ 怎麼換，見表後的說明。

\begin{center}
\begin{tabular}{p{0.15\linewidth}p{0.33\linewidth}p{0.42\linewidth}}
\toprule
項目 & 虛擬碼 & C \\
\midrule
陣列編號 & 從 $1$ 開始：$A[1..n]$ & 從 $0$ 開始：\texttt{a[0..n-1]} \\
指定 & $x \gets x + 1$ & \texttt{x = x + 1} \\
比較相等 & $x = y$、不等 $x \ne y$ & \texttt{x == y}、\texttt{x != y} \\
計數迴圈 & \textbf{for} $i \gets 1$ \textbf{to} $n$（包含 $n$）&
\texttt{for (int i = 1; i <= n; i++)} \\
往下數 & \textbf{for} $i \gets n$ \textbf{downto} $1$（包含 $1$）&
\texttt{for (int i = n; i >= 1; i--)} \\
且、或 & \textbf{and}、\textbf{or}，左邊決定結果時右邊不算 & \texttt{\&\&}、\texttt{||}，
一樣不算右邊 \\
空指標 & \NIL & \texttt{NULL} \\
物件的欄位 & $x.\mathit{key}$，$x$ 是指向物件的指標 & \texttt{x->key} \\
程序名稱 & \proc{Insertion-Sort}$(A, n)$，名稱用小型大寫 & \texttt{insertion\_sort(a, n)} \\
回傳 & \textbf{return} $x$ & \texttt{return x;} \\
註解 & 寫在行尾，以 // 開頭 & \texttt{//} 或 \texttt{/* */} \\
區塊 & 用縮排表示，不寫 \textbf{end} & 用大括號 \\
\bottomrule
\end{tabular}
\end{center}

另外幾個約定：

\begin{itemize}
\item 虛擬碼的變數都是區域變數；陣列和物件傳進程序時，傳的是指標，程序裡改了內
容，呼叫的人看得到。 C 也一樣，陣列參數其實是指標。
\item 虛擬碼裡的數學運算不會溢位，除法要寫明是 $\lfloor a / b \rfloor$ 還是
$\lceil a / b \rceil$。 C 程式要自己注意溢位，會超過 \texttt{int} 的計數一律用
\texttt{long long}。
\item 虛擬碼可以寫一句中文代替一段細節，例如「建立陣列 $\mathit{seen}[0..M-1]$，
每一格都設成 \FALSE」，只要那一句的成本清楚。
\item 好幾個字母的變數名稱用斜體整個寫出，例如 $\mathit{key}$、$\mathit{seen}$，指的是一
個變數，不是幾個字母相乘。
\item $\lg n$ 表示 $\log_2 n$，$\ln n$ 表示自然對數。
\end{itemize}

把虛擬碼翻成 C 時最容易錯的地方是編號差 $1$。 這本書的做法是：C 程式的陣列一律從
$0$ 開始，虛擬碼的 $A[i]$ 對應 \texttt{a[i - 1]}，翻譯時整段一起改，不要一行一行
加減 $1$。 例如虛擬碼的 \textbf{for} $i \gets 2$ \textbf{to} $n$ 處理第 $2$ 到第
$n$ 個元素，翻成 C 是 \texttt{for (i = 1; i < n; i++)}，迴圈裡的 $A[i]$ 直接寫
\texttt{a[i]}。

\section{插入排序：逐行計算成本}
\label{sec:ana-insertion}

\textbf{排序}（sorting）的問題是：輸入 $n$ 個數 $a_1, a_2, \ldots, a_n$，輸出它們由
小到大的排列。 第~14~章和第~15~章會講十種左右的排序方法；這一節先用最簡單的一種，示
範怎麼逐行算出執行次數。

\textbf{插入排序}（insertion sort）的想法是讓陣列左邊有一段已經排好。 一開始這段
只有第一個數。 每一輪把這段右邊的下一個數拿出來，叫它 $\mathit{key}$，從這段的最
右邊往左看，比 $\mathit{key}$ 大的數都往右挪一格，直到遇到不比它大的數，或是走到
最左邊，再把 $\mathit{key}$ 放進空出來的位置。 這樣排好的一段就多了一個數。

\begin{pseudo}[label=alg:insertion]{插入排序}
\Procedure{Insertion-Sort}{$A, n$}
    \For{$i \gets 2$ \To $n$}
        \State $\mathit{key} \gets A[i]$
        \State $j \gets i - 1$
        \While{$j \ge 1$ \AND $A[j] > \mathit{key}$}
            \State $A[j + 1] \gets A[j]$ \Comment{往右挪一格}
            \State $j \gets j - 1$
        \EndWhile
        \State $A[j + 1] \gets \mathit{key}$
    \EndFor
\EndProcedure
\end{pseudo}

\subsection{手算追蹤}

用 $A = \langle 27, 12, 35, 6, 19, 8 \rangle$ 跑一次。 下圖每一列是外層迴圈做完一
輪之後的陣列，藍底是已經排好的一段，橘框是這一輪插進去的 $\mathit{key}$。

\begin{center}
\begin{tikzpicture}
\insrow{0}{開始}{27,12,35,6,19,8}{1}{0}
\insrow{-1.1}{$i = 2$，$\mathit{key} = 12$}{12,27,35,6,19,8}{2}{1}
\insrow{-2.2}{$i = 3$，$\mathit{key} = 35$}{12,27,35,6,19,8}{3}{3}
\insrow{-3.3}{$i = 4$，$\mathit{key} = 6$}{6,12,27,35,19,8}{4}{1}
\insrow{-4.4}{$i = 5$，$\mathit{key} = 19$}{6,12,19,27,35,8}{5}{3}
\insrow{-5.5}{$i = 6$，$\mathit{key} = 8$}{6,8,12,19,27,35}{6}{2}
\foreach \k in {1,...,6} { \node[idxlbl] at (\k*0.9, -6.2) {\k}; }
\end{tikzpicture}
\diagcaption{插入排序每一輪之後的陣列。 藍底是已排好的前段，橘框是這一輪的
$\mathit{key}$ 最後放的位置。}
\label{fig:ana-insertion-trace}
\end{center}

每一輪的細節列成下表。 「比較」是第 5 行 $A[j] > \mathit{key}$ 實際比了幾次（$j$
已經變成 $0$ 時，因為 \textbf{and} 左邊已經不成立，右邊不算），「挪動」是第 6 行
執行的次數。

\begin{center}
\begin{tabular}{cccccl}
\toprule
$i$ & $\mathit{key}$ & 比較 & 挪動 & 停下的原因 & 這一輪之後 \\
\midrule
2 & 12 & 1 & 1 & $j = 0$ & 12 27 35 6 19 8 \\
3 & 35 & 1 & 0 & $A[2] = 27 \le 35$ & 12 27 35 6 19 8 \\
4 & 6 & 3 & 3 & $j = 0$ & 6 12 27 35 19 8 \\
5 & 19 & 3 & 2 & $A[2] = 12 \le 19$ & 6 12 19 27 35 8 \\
6 & 8 & 5 & 4 & $A[1] = 6 \le 8$ & 6 8 12 19 27 35 \\
\midrule
合計 & & 13 & 10 & & \\
\bottomrule
\end{tabular}
\end{center}

\begin{examnote}
考卷上的追蹤題，通常要求寫出「每一輪之後」的陣列，像上表最右欄那樣一輪一列。 先
看清楚題目要的是每一輪之後，還是每一次交換、挪動之後，兩者的列數不同。 題目問
比較次數時，記得 $j$ 走到 $0$ 那一次不算比較。

有的課本把 $\mathit{key}$ 先放進 $A[0]$ 當\textbf{哨兵}（sentinel），迴圈條件只寫
$A[j] > \mathit{key}$，省掉 $j \ge 1$ 的檢查。 這個版本裡 $j$ 走到 $0$ 時會和哨兵比
一次，所以第 $i$ 輪最多比 $i$ 次；上面的例子照這個版本數，比較次數是 $15$，不是
$13$。 題目給了虛擬碼時，照題目的版本數。
\end{examnote}

\subsection{C 程式}

下面的程式從輸入讀進好幾組陣列，每組排序之後印出結果，以及比較和挪動的次數。
C 的陣列從 $0$ 開始，所以外層從 \texttt{i = 1} 開始，內層的條件 $j \ge 1$ 變成
\texttt{j >= 0}。 為了數比較次數，把 \textbf{and} 的兩半拆開：\texttt{j >= 0}
寫在 \texttt{while} 的條件，\texttt{a[j] <= key} 寫在迴圈裡面，比之前先加一。

\begin{example}[label=ex:insertion-count]{c}{(insertion\_count.c) 插入排序並計算比較和挪動的次數}
#include <stdio.h>
#include <stdlib.h>

// 插入排序，a[0..n-1] 由小到大。*cmp 記錄 a[j] > key 比較了幾次，
// *shift 記錄往右搬了幾次。
void insertion_sort(int *a, int n, long long *cmp, long long *shift)
{
    *cmp = 0;
    *shift = 0;
    for (int i = 1; i < n; i++) {
        int key = a[i];
        int j = i - 1;
        while (j >= 0) {
            (*cmp)++;
            if (a[j] <= key) {
                break;
            }
            a[j + 1] = a[j];
            (*shift)++;
            j--;
        }
        a[j + 1] = key;
    }
}

int main(void)
{
    int t;
    if (scanf("%d", &t) != 1) {
        return 1;
    }
    while (t--) {
        int n;
        if (scanf("%d", &n) != 1 || n <= 0) {
            return 1;
        }
        int *a = malloc(n * sizeof *a);
        if (a == NULL) {
            return 1;
        }
        for (int i = 0; i < n; i++) {
            if (scanf("%d", &a[i]) != 1) {
                free(a);
                return 1;
            }
        }
        long long cmp, shift;
        insertion_sort(a, n, &cmp, &shift);
        for (int i = 0; i < n; i++) {
            printf("%d%c", a[i], i + 1 < n ? ' ' : '\n');
        }
        printf("comparisons = %lld, shifts = %lld\n", cmp, shift);
        free(a);
    }
    return 0;
}
\end{example}
\iosample{3 \\ 6 \\ 4 9 13 21 30 38 \\ 6 \\ 38 30 21 13 9 4 \\ 6 \\ 27 12 35 6 19 8}{4 9 13 21 30 38 \\ comparisons = 5, shifts = 0 \\ 4 9 13 21 30 38 \\ comparisons = 15, shifts = 15 \\ 6 8 12 19 27 35 \\ comparisons = 13, shifts = 10}

三組輸入都是 $6$ 個數，比較次數卻從 $5$ 到 $15$。 已經排好的陣列每一輪只比一次就
停；完全相反的陣列每一輪都要一路挪到最左邊。 執行時間不只和 $n$ 有關，也和輸入
長什麼樣子有關，這一點在第~\ref{sec:ana-cases}~節會再整理。

\subsection{逐行計算}
\label{sec:ana-line-cost}

現在對一般的 $n$ 算出虛擬碼~\ref{alg:insertion} 的執行時間。 假設第 $k$ 行執行一
次花 $c_k$ 的時間（$c_k$ 是常數，第 2 行到第 8 行各自不同），再數每一行執行幾次。
外層的 $i$ 跑 $2$ 到 $n$。 對固定的 $i$，令 $t_i$ 是第 5 行 \textbf{while} 的條件
被檢查的次數：挪了幾次就檢查幾次，再加上最後讓迴圈停下來的那一次，所以挪動
$t_i - 1$ 次。

\begin{center}
\begin{tabular}{clcc}
\toprule
行 & 內容 & 每次的成本 & 執行次數 \\
\midrule
2 & \textbf{for} $i \gets 2$ \textbf{to} $n$ & $c_2$ & $n$ \\
3 & $\mathit{key} \gets A[i]$ & $c_3$ & $n - 1$ \\
4 & $j \gets i - 1$ & $c_4$ & $n - 1$ \\
5 & \textbf{while} $j \ge 1$ \textbf{and} $A[j] > \mathit{key}$ & $c_5$ &
$\sum_{i=2}^{n} t_i$ \\
6 & $A[j + 1] \gets A[j]$ & $c_6$ & $\sum_{i=2}^{n} (t_i - 1)$ \\
7 & $j \gets j - 1$ & $c_7$ & $\sum_{i=2}^{n} (t_i - 1)$ \\
8 & $A[j + 1] \gets \mathit{key}$ & $c_8$ & $n - 1$ \\
\bottomrule
\end{tabular}
\end{center}

第 2 行執行 $n$ 次：$i = 2, 3, \ldots, n$ 各檢查一次，再加上 $i = n + 1$ 時發現超過
而離開的那一次。 把每一行的成本乘上次數加起來，得到
\[
T(n) = c_2 n + (c_3 + c_4 + c_8)(n - 1) + c_5 \sum_{i=2}^{n} t_i
+ (c_6 + c_7) \sum_{i=2}^{n} (t_i - 1).
\]

\textbf{最好的情況}：陣列本來就排好了。 每一輪第一次檢查就發現 $A[i-1] \le
\mathit{key}$，$t_i = 1$，挪動 $0$ 次。
\[
T(n) = c_2 n + (c_3 + c_4 + c_5 + c_8)(n - 1) = (c_2 + c_3 + c_4 + c_5 + c_8)\, n
- (c_3 + c_4 + c_5 + c_8).
\]
這是 $an + b$ 的形式，$a$、$b$ 是由各個 $c_k$ 決定的常數，執行時間是 $n$ 的一次
函數。

\textbf{最壞的情況}：陣列由大到小排。 第 $i$ 輪的 $\mathit{key}$ 比前面每一個數都
小，要挪 $i - 1$ 次，最後在 $j = 0$ 時停下，$t_i = i$。 用
\[
\sum_{i=2}^{n} i = \frac{n(n+1)}{2} - 1, \qquad
\sum_{i=2}^{n} (i - 1) = \frac{n(n-1)}{2}
\]
代進去：
\[
T(n) = \left(\frac{c_5 + c_6 + c_7}{2}\right) n^2
+ \left(c_2 + c_3 + c_4 + c_8 + \frac{c_5 - c_6 - c_7}{2}\right) n
- (c_3 + c_4 + c_5 + c_8).
\]
這是 $an^2 + bn + c$ 的形式，執行時間是 $n$ 的二次函數。

這些 $c_k$ 是多少，要看機器和編譯器，算不出來，也不重要。 $n$ 很大時，$an^2$ 這一
項會遠遠超過其他項，係數 $a$ 只影響「大多少倍」，不影響「$n$ 變兩倍時時間變幾
倍」。 下一節介紹的漸近符號，就是把這種「只看最高次項、不看係數」的說法寫得精確。

\subsection{挪動次數和逆序數}

在陣列裡，如果 $i < j$ 而 $A[i] > A[j]$，就說 $(i, j)$ 是一個\textbf{逆序對}
（inversion）。 插入排序每挪動一次，剛好消掉一個逆序對：$A[j]$ 往右越過
$\mathit{key}$，原本 $A[j]$ 在 $\mathit{key}$ 左邊而且比較大，挪完之後兩者的順序對
了，而其他數對的前後關係都沒有變。 排好的陣列沒有逆序對，所以
\[
\text{插入排序的挪動次數} = \text{輸入陣列的逆序對個數}.
\]
例如 $\langle 27, 12, 35, 6, 19, 8 \rangle$ 的逆序對有 $(27,12)$、$(27,6)$、
$(27,19)$、$(27,8)$、$(12,6)$、$(12,8)$、$(35,6)$、$(35,19)$、$(35,8)$、$(19,8)$，
共 $10$ 個（這裡直接寫數值，不寫位置），和程式印出的 \texttt{shifts = 10} 一樣。
$n$ 個數最多有 $\binom{n}{2} = n(n-1)/2$ 個逆序對，發生在完全相反的時候。

\section{迴圈不變量}
\label{sec:ana-invariant}

算完執行時間，還要說明插入排序真的會把陣列排好。 測幾組輸入只能說明那幾組是對
的；要說明每一組都對，常用的工具是\textbf{迴圈不變量}（loop invariant）：一個和迴
圈變數有關的敘述，每一輪開始時都成立。 證明分成三步。

\begin{enumerate}
\item \textbf{初始}：第一輪開始之前，敘述成立。
\item \textbf{維持}：如果某一輪開始時敘述成立，那麼這一輪做完、下一輪開始時，敘述
仍然成立。
\item \textbf{結束}：迴圈停下來時，把敘述和「讓迴圈停下來的條件」合在一起，可以推
出我們要的結論。
\end{enumerate}

前兩步和數學歸納法一樣，第三步是歸納法沒有的：歸納法證明對所有正整數都成立，迴
圈卻會停，停的時候要把不變量用起來。

\textbf{插入排序的不變量}：外層迴圈每一輪開始時（$i$ 是這一輪的值），$A[1..i-1]$
裡的數就是原本放在 $A[1..i-1]$ 的那些數，而且已經由小到大排好。

\begin{itemize}
\item \textbf{初始}：第一輪 $i = 2$，$A[1..1]$ 只有一個數，就是原本的 $A[1]$，一
個數一定排好了。
\item \textbf{維持}：這一輪把 $\mathit{key} = A[i]$ 存起來，內層迴圈把
$A[i-1], A[i-2], \ldots$ 之中比 $\mathit{key}$ 大的數依序往右挪一格，停在第一個
不比 $\mathit{key}$ 大的 $A[j]$（或 $j = 0$）。 因為 $A[1..i-1]$ 原本排好了，被
挪的是最右邊連續的一段，也就是所有比 $\mathit{key}$ 大的數；它們挪動後順序不
變。 $\mathit{key}$ 放進 $A[j+1]$ 之後，左邊的數都不大於它，右邊的數都大於它，
所以 $A[1..i]$ 排好了，而且裡面的數正是原本 $A[1..i]$ 的那些數。 下一輪的 $i$
加一，不變量成立。（嚴格的寫法要對內層迴圈另外寫一個不變量，這裡用文字帶過。）
\item \textbf{結束}：外層迴圈在 $i = n + 1$ 時停下。 代進不變量：$A[1..n]$ 是原本
的 $n$ 個數，而且排好了，這正是排序要的結果。
\end{itemize}

\begin{examnote}
用迴圈不變量證明時，三步都要寫出來，而且不變量要寫成和迴圈變數有關的具體敘述
（「$A[1..i-1]$ 排好了」），不能只寫「陣列越來越有序」。 結束那一步要明白寫出迴
圈停下時變數的值（這裡是 $i = n + 1$），再代進不變量。
\end{examnote}

不變量要選得剛好：太弱推不出結論，太強維持不住。 例如只寫「$A[1..i-1]$ 排好了」，
少了「是原本的那些數」，那麼一個把整個陣列都改成 $0$ 的程式也滿足它。
觀念題第~\ref{ana-q-lower-bound}~題的二分搜尋、第~14~章快速排序的分割（partition）都會再用到這個工具。

\section{漸近符號}
\label{sec:ana-notation}

第~\ref{sec:ana-line-cost}~節算出插入排序最壞的情況是 $an^2 + bn + c$。 $n = 10^4$
時，$an^2$ 比 $bn$ 大了大約 $10^4 \cdot a/b$ 倍，只要 $a/b$ 不是非常小，低次項
幾乎沒有影響。 所以我們只想說「它是 $n^2$ 這一級的」。 這一節把這句話定義清楚。

這一節的函數都定義在正整數 $n$ 上，而且 $n$ 夠大以後取值不是負的（執行時間不會是
負的）。

\subsection{$\Theta$、$O$、$\Omega$}

\begin{definition}[$\Theta$]
\label{def:theta}
$\Theta(g(n))$ 是所有滿足下面條件的函數 $f(n)$ 組成的集合：存在正的常數 $c_1$、
$c_2$ 和 $n_0$，使得所有 $n \ge n_0$ 都有
\[
0 \le c_1 g(n) \le f(n) \le c_2 g(n).
\]
\end{definition}

白話是：$n$ 夠大以後，$f(n)$ 被夾在 $g(n)$ 的兩個常數倍之間。 $n_0$ 以前怎麼樣都
不管。 $f(n) \in \Theta(g(n))$ 習慣寫成 $f(n) = \Theta(g(n))$，這裡的等號意思是
「屬於」，不能左右對調：寫 $\Theta(g(n)) = f(n)$ 沒有意義。

\begin{center}
\begin{tikzpicture}[x=0.85cm, y=0.7cm]
\draw[-{Stealth[length=2mm]}] (0,0) -- (8.6,0) node[right, font=\small] {$n$};
\draw[-{Stealth[length=2mm]}] (0,0) -- (0,8.2);
\draw[thick, coreStroke, domain=0:8, samples=40] plot (\x, {1.4*(0.08*\x*\x+0.5)})
    node[right, font=\small] {$c_2 g(n)$};
\draw[thick, coreStroke, dashed, domain=0:8, samples=40] plot (\x, {0.6*(0.08*\x*\x+0.5)})
    node[right, font=\small] {$c_1 g(n)$};
\draw[thick, memStroke] plot[smooth] coordinates {(0.2,2.0) (1,1.6) (2,0.95) (3,1.3)
    (4,1.9) (5,2.3) (6,3.6) (7,4.2) (8,5.6)} node[right, font=\small] {$f(n)$};
\draw[densely dotted] (2,0) -- (2,7.6);
\node[below, font=\small] at (2,0) {$n_0$};
\end{tikzpicture}
\diagcaption{$f(n) = \Theta(g(n))$：在 $n_0$ 右邊，$f(n)$ 一直落在 $c_1 g(n)$ 和
$c_2 g(n)$ 之間。 $n_0$ 左邊的部分不影響結論。}
\label{fig:ana-theta}
\end{center}

只要上界或只要下界時，分別用 $O$ 和 $\Omega$。

\begin{definition}[$O$ 和 $\Omega$]
\label{def:big-o}
$O(g(n))$ 是所有滿足下面條件的 $f(n)$：存在正的常數 $c$ 和 $n_0$，使得所有
$n \ge n_0$ 都有 $0 \le f(n) \le c\, g(n)$。

$\Omega(g(n))$ 是所有滿足下面條件的 $f(n)$：存在正的常數 $c$ 和 $n_0$，使得所有
$n \ge n_0$ 都有 $0 \le c\, g(n) \le f(n)$。
\end{definition}

直接從定義可以得到：

\begin{theorem}
\label{thm:theta-iff}
$f(n) = \Theta(g(n))$ 的充分必要條件是 $f(n) = O(g(n))$ 而且 $f(n) = \Omega(g(n))$。
\end{theorem}

\begin{proof}
如果 $f(n) = \Theta(g(n))$，定義裡的 $c_2$、$n_0$ 就是 $O$ 要的常數，$c_1$、$n_0$
就是 $\Omega$ 要的常數。 反過來，假設 $O$ 的常數是 $c$、$n_1$，$\Omega$ 的常數是
$c'$、$n_2$，取 $c_1 = c'$、$c_2 = c$、$n_0 = \max(n_1, n_2)$，$n \ge n_0$ 時兩個
不等式同時成立。
\end{proof}

\subsection{用定義證明}

\textbf{例 1}：證明 $4n^2 + 7n - 30 = \Theta(n^2)$。

上界：$n \ge 1$ 時 $7n \le 7n^2$，而 $-30 < 0$，所以
$4n^2 + 7n - 30 \le 4n^2 + 7n^2 = 11n^2$。 取 $c_2 = 11$。

下界：想要 $4n^2 + 7n - 30 \ge 3n^2$，也就是 $n^2 + 7n \ge 30$。 $n = 3$ 時左邊是
$30$，剛好成立，$n$ 更大時左邊更大，所以 $n \ge 3$ 時成立。 取 $c_1 = 3$。

兩邊都要成立，取 $n_0 = 3$。 另外 $n \ge 3$ 時 $3n^2 \ge 0$，定義裡的 $0 \le$ 也滿
足。 所以 $c_1 = 3$、$c_2 = 11$、$n_0 = 3$ 符合定義~\ref{def:theta}。

常數的選法不只一組。 下界也可以取 $c_1 = 1$，這時需要 $3n^2 + 7n \ge 30$，
$n = 2$ 時左邊只有 $26$，要到 $n \ge 3$ 才成立。 證明只要找到一組，並且說清楚每個不等式從哪一個 $n$ 開始成立。

\begin{examnote}
用定義證明 $\Theta$、$O$、$\Omega$ 時，答案裡要明白寫出 $c_1$、$c_2$（或 $c$）
和 $n_0$ 的數值，並驗證「所有 $n \ge n_0$」都成立，不能只代一兩個 $n$ 檢查。 常
用的手法是把負的項丟掉得到上界，把低次項併進最高次項的一部分得到下界。
\end{examnote}

\textbf{例 2}：證明 $n^3 \ne O(n^2)$。

用反證法。 假設有正的 $c$、$n_0$，使得所有 $n \ge n_0$ 都有 $n^3 \le c\,n^2$，兩
邊除以 $n^2$ 得到 $n \le c$。 但取 $n = \max(n_0, \lceil c \rceil + 1)$，這個 $n$
大於等於 $n_0$，卻比 $c$ 大，矛盾。

要證明「不是 $O$」，就要說明不管常數怎麼選都不行，所以通常用反證法：假設常數存
在，再找一個夠大的 $n$ 推出矛盾。

\textbf{例 3}：多項式只看最高次。 設
$p(n) = a_d n^d + a_{d-1} n^{d-1} + \cdots + a_0$，$a_d > 0$。 令
$A = |a_{d-1}| + \cdots + |a_0|$。 $n \ge 1$ 時每個 $|a_k| n^k \le |a_k| n^{d-1}$，所以
\[
a_d n^d - A n^{d-1} \le p(n) \le a_d n^d + A n^{d-1}.
\]
右邊不超過 $(a_d + A) n^d$；左邊在 $n \ge 2A / a_d$ 時至少是 $\frac{a_d}{2} n^d$。
所以 $p(n) = \Theta(n^d)$。 插入排序最壞情況的 $an^2 + bn + c$（$a > 0$）是
$\Theta(n^2)$，最好情況的 $an + b$ 是 $\Theta(n)$。

\subsection{$O$ 只是上界}

$O$ 只說「不會比這個大太多」，沒有說「差不多這麼大」。 $5n = O(n^2)$、
$5n = O(n^3)$ 都對，因為 $n \ge 5$ 時 $5n \le n^2 \le n^3$。 所以「插入排序的執行
時間是 $O(n^3)$」是一句對的話，只是不夠精確。

考卷上的選擇題常常利用這一點。 讀到 $O$ 時，先問自己：這句話說的是上界，題目給
的函數有沒有超過這個上界？ 不要把 $O$ 讀成「剛好是」。 想說「剛好是這一級」時用
$\Theta$。

\subsection{$o$ 和 $\omega$}

$O$ 允許 $f$ 和 $g$ 同一級，例如 $2n^2 = O(n^2)$。 想說「$f$ 嚴格比 $g$ 小一級」時，
用小寫的 $o$。

\begin{definition}[$o$ 和 $\omega$]
\label{def:little-o}
$o(g(n))$ 是所有滿足下面條件的 $f(n)$：\textbf{對每一個}正的常數 $c$，都存在
$n_0$，使得所有 $n \ge n_0$ 都有 $0 \le f(n) < c\, g(n)$。

$\omega(g(n))$ 是所有滿足下面條件的 $f(n)$：\textbf{對每一個}正的常數 $c$，都存在
$n_0$，使得所有 $n \ge n_0$ 都有 $0 \le c\, g(n) < f(n)$。
\end{definition}

和 $O$ 的差別在「存在一個 $c$」換成了「每一個 $c$」：$c$ 再小，$n$ 夠大以後
$f(n)$ 都比 $c\,g(n)$ 小。 例如 $2n = o(n^2)$：給定 $c$，取 $n_0 > 2/c$，$n \ge n_0$
時 $2n < c\,n^2$。 但 $2n^2 \ne o(n^2)$：取 $c = 1$，$2n^2 < n^2$ 永遠不成立。

用極限判斷常常比較快。

\begin{theorem}
\label{thm:limit}
設 $g(n) > 0$，而且極限 $L = \lim_{n \to \infty} f(n)/g(n)$ 存在（可以是 $\infty$）。
\begin{enumerate}
\item $L = 0$：$f(n) = o(g(n))$，因此 $f(n) = O(g(n))$，但 $f(n) \ne \Omega(g(n))$。
\item $0 < L < \infty$：$f(n) = \Theta(g(n))$。
\item $L = \infty$：$f(n) = \omega(g(n))$，因此 $f(n) = \Omega(g(n))$，但 $f(n) \ne O(g(n))$。
\end{enumerate}
\end{theorem}

第 2 點的理由：$n$ 夠大以後 $f(n)/g(n)$ 落在 $[L/2, 2L]$ 之間，取 $c_1 = L/2$、
$c_2 = 2L$。 第 1 點：極限是 $0$ 的意思是，對每一個 $c > 0$，$n$ 夠大以後
$f(n)/g(n) < c$，這正是 $o$ 的定義；取 $c = 1$ 就得到 $O$。 如果又有
$f(n) \ge c' g(n)$（$c' > 0$），比值會一直不小於 $c'$，不可能趨近 $0$，所以
$f(n) \ne \Omega(g(n))$。 第 3 點的理由相同，大小對調。

要注意這個定理只有一個方向：極限不存在時，什麼都不能說，要回到定義。 例如
\[
f(n) = \begin{cases} n^2, & n \text{ 是偶數} \\ n, & n \text{ 是奇數} \end{cases}
\]
和 $g(n) = n^2$ 相比，$f(n)/g(n)$ 在 $1$ 和 $1/n$ 之間跳，極限不存在。 但 $f(n) \le
n^2$ 對所有 $n$ 成立，所以 $f(n) = O(n^2)$；奇數的 $n$ 讓 $f(n) = n$，所以 $f(n) \ne
\Omega(n^2)$。

\subsection{性質}

下面幾條性質直接從定義推得，之後會常常用到。 $f$、$g$、$h$ 都是夠大以後非負的函數。

\begin{itemize}
\item \textbf{遞移}：$f = O(g)$ 而且 $g = O(h)$，則 $f = O(h)$。 $\Theta$、$\Omega$、
$o$、$\omega$ 也一樣。
\item \textbf{自反}：$f = \Theta(f)$，也就是 $f = O(f)$、$f = \Omega(f)$。
\item \textbf{對稱}：$f = \Theta(g)$ 的充分必要條件是 $g = \Theta(f)$。
\item \textbf{轉置}：$f = O(g)$ 的充分必要條件是 $g = \Omega(f)$；$f = o(g)$ 的充分
必要條件是 $g = \omega(f)$。
\item \textbf{加法}：$f + g = \Theta(\max(f, g))$。 例如兩段程式接在一起，總時間
由比較慢的那一段決定。
\end{itemize}

這些性質很像數字的大小關係（$O$ 像 $\le$，$\Theta$ 像 $=$，$o$ 像 $<$），但有一點
不一樣：任兩個數一定比得出大小，兩個函數卻不一定。 例如 $g(n) = n^2$，而
$f(n)$ 在偶數時是 $n$、奇數時是 $n^3$。 偶數的 $n$ 讓 $f$ 遠小於 $g$，奇數的 $n$
讓 $f$ 遠大於 $g$，所以 $f \ne O(g)$，$f \ne \Omega(g)$。

\subsection{等式裡的漸近符號}

漸近符號出現在式子的一部分時，代表「某一個屬於這個集合的函數」。 例如
\[
T(n) = 2n^2 + \Theta(n)
\]
是說 $T(n) = 2n^2 + h(n)$，其中 $h(n)$ 是某個 $\Theta(n)$ 的函數。 這種寫法可以保
留關心的部分（最高次的係數 $2$），把不關心的細節收起來。 第~2~章寫遞迴式
$T(n) = 2T(n/2) + \Theta(n)$ 也是這個意思。

式子兩邊都有漸近符號時，例如 $2n^2 + \Theta(n) = \Theta(n^2)$，意思是：左邊不管
$\Theta(n)$ 換成哪一個函數，右邊都找得到一個 $\Theta(n^2)$ 的函數讓等式成立。 這仍
然是單向的，右邊比左邊粗。

\section{常見函數的成長速度}
\label{sec:ana-functions}

分析演算法時，最常出現的函數是對數、多項式、指數和階乘。 這一節整理它們之間的
大小關係，最後給一張由慢到快的表。

\subsection{對數}

$\lg n = \log_2 n$。 換底公式 $\log_a n = \log_b n / \log_b a$ 說明不同底的對數只差
一個常數倍，所以在 $O$、$\Theta$ 裡面寫 $\log n$ 不用寫底：$\log_{10} n =
\Theta(\lg n) = \Theta(\ln n)$。

底數在指數上就不能省了。 $4^n = 2^{2n} = (2^n)^2$，比 $2^n$ 大了 $2^n$ 倍，所以
$4^n \ne O(2^n)$。 另外幾個常用的換算：
\[
2^{\lg n} = n, \qquad 4^{\lg n} = n^2, \qquad a^{\log_b n} = n^{\log_b a}.
\]
最後一個等式兩邊取 $\log_b$ 就看得出來，第~2~章的主定理會一直用到它。

\subsection{對數、多項式、指數}

三條常用的結論，$a$、$b$、$\varepsilon$ 是常數：

\begin{enumerate}
\item 任何次方的對數都比任何正次方的多項式慢：$(\lg n)^a = o(n^{\varepsilon})$，
$a > 0$，$\varepsilon > 0$。
\item 任何多項式都比任何底大於 $1$ 的指數慢：$n^a = o(b^n)$，$b > 1$。
\item 多項式之間只看次方：$a < b$ 時 $n^a = o(n^b)$。
\end{enumerate}

第 1 點可以用定理~\ref{thm:limit}：令 $n = 2^m$，$(\lg n)^a / n^{\varepsilon} =
m^a / 2^{\varepsilon m}$，這是第 2 點的情形（以 $m$ 為變數，底是 $2^{\varepsilon}$），
所以極限是 $0$。 第 2 點可以對 $n^a / b^n$ 用羅必達法則 $\lceil a \rceil$ 次。

這幾條結論的數量感要記住。 $(\lg n)^3$ 在 $n = 10^6$ 時大約是 $7900$，比
$\sqrt{n} = 1000$ 還大；可是 $n = 2^{100}$ 時，$(\lg n)^3 = 10^6$，$\sqrt n = 2^{50}
\approx 10^{15}$。 漸近的比較說的是 $n$ 夠大以後，「夠大」可能非常大。

\subsection{階乘}

$n!$ 比任何 $b^n$ 都快，又比 $n^n$ 慢。 比較精確的估計是\textbf{斯特林公式}
（Stirling's approximation）：
\[
n! = \sqrt{2\pi n}\left(\frac{n}{e}\right)^n \left(1 + \Theta\!\left(\frac{1}{n}\right)\right).
\]
排序的下界（第~15~章）要用到 $\lg(n!)$。 它其實不需要斯特林公式也能估：
\[
n! = 1 \cdot 2 \cdots n \le n^n,
\qquad
n! \ge \left\lceil \tfrac{n}{2} \right\rceil \cdot \left(\left\lceil \tfrac{n}{2} \right\rceil + 1\right) \cdots n
\ge \left(\frac{n}{2}\right)^{n/2}.
\]
第二個式子只留下後半段的乘數，每個都至少 $n/2$，個數至少 $n/2$。 兩邊取 $\lg$：
\[
\frac{n}{2} \lg \frac{n}{2} \le \lg(n!) \le n \lg n,
\]
左邊在 $n \ge 4$ 時至少是 $\frac{n}{4}\lg n$（因為 $\lg(n/2) = \lg n - 1 \ge
\frac{1}{2}\lg n$）。 所以 $\lg(n!) = \Theta(n \lg n)$。

\subsection{級數}

逐行分析時最常出現下面幾種和。

\begin{itemize}
\item \textbf{等差}：$\displaystyle\sum_{i=1}^{n} i = \frac{n(n+1)}{2} = \Theta(n^2)$。 一
般地，$\displaystyle\sum_{i=1}^{n} i^k = \Theta(n^{k+1})$，$k \ge 0$ 是常數。
\item \textbf{等比}：$r \ne 1$ 時
$\displaystyle\sum_{i=0}^{m} r^i = \frac{r^{m+1} - 1}{r - 1}$。 $0 < r < 1$ 時，不管
$m$ 多大，和都小於 $\frac{1}{1-r}$，是 $\Theta(1)$；$r > 1$ 時，和是
$\Theta(r^m)$，由最後一項決定。
\item \textbf{調和}：$H_n = \displaystyle\sum_{i=1}^{n} \frac{1}{i}$。 拿 $1/x$ 的
積分上下夾，得到 $\ln(n+1) \le H_n \le \ln n + 1$，所以 $H_n = \Theta(\lg n)$。
\end{itemize}

等比級數的兩個結論要特別記住：公比小於 $1$ 時，總和和第一項同一級；公比大於 $1$ 時，
總和和最後一項同一級。 下一節分析「每次減半」「每次加倍」的迴圈，都是用這兩句話。

\subsection{由慢到快}

圖~\ref{fig:ana-growth} 畫出 $n \le 20$ 時幾個常見函數的大小，下表再列出 $n$ 更大時
的值。 在一般的電腦上，一秒大約能做 $10^8$ 到 $10^9$ 個基本運算。 對照下表，
$n = 10^6$ 時 $n \lg n$ 的演算法一瞬間就跑完，$n^2$ 要十幾分鐘到幾個小時，指數的演
算法連 $n = 100$ 都跑不完。 演算法課關心的，主要就是能不能從表的下半部換到上半部。

\begin{center}
\begin{tikzpicture}[x=0.42cm, y=0.075cm]
\begin{scope}
\clip (0,0) rectangle (21,62);
\draw[growth, black!60, domain=1:20, samples=60] plot (\x, {ln(\x)/ln(2)});
\draw[growth, coreStroke, domain=1:20, samples=60] plot (\x, {\x});
\draw[growth, greenStroke, domain=1:20, samples=60] plot (\x, {\x*ln(\x)/ln(2)});
\draw[growth, memStroke, domain=1:20, samples=60] plot (\x, {\x*\x});
\draw[growth, optStroke, domain=1:6, samples=60] plot (\x, {pow(2,\x)});
\end{scope}
\draw[-{Stealth[length=2mm]}] (0,0) -- (21.5,0) node[right, font=\small] {$n$};
\draw[-{Stealth[length=2mm]}] (0,0) -- (0,64);
\foreach \t in {5,10,15,20} { \draw (\t,0) -- (\t,-1) node[below, font=\scriptsize] {\t}; }
\foreach \t in {20,40,60} { \draw (0,\t) -- (-0.3,\t) node[left, font=\scriptsize] {\t}; }
\node[glabel, text=black!60] at (20.1,4.3) {$\lg n$};
\node[glabel, text=coreStroke] at (20.1,20) {$n$};
\node[glabel, text=greenStroke] at (15.4,60) {$n \lg n$};
\node[glabel, text=memStroke] at (7.9,60) {$n^2$};
\node[glabel, text=optStroke] at (6.0,60) {$2^n$};
\end{tikzpicture}
\diagcaption{$n \le 20$ 時幾個常見函數的大小（縱軸只畫到 $60$）。 $2^n$ 在
$n = 6$ 左右就超出圖外，$n^2$ 在 $n = 8$ 左右，$n \lg n$ 在 $n = 15$ 左右。}
\label{fig:ana-growth}
\end{center}

\begin{center}
\begin{tabular}{lrrrl}
\toprule
函數 & $n = 10$ & $n = 10^3$ & $n = 10^6$ & 常見的來源 \\
\midrule
$1$ & 1 & 1 & 1 & 陣列取值、堆疊的推入 \\
$\lg n$ & 3.3 & 10.0 & 19.9 & 二分搜尋、平衡樹的一次操作 \\
$\sqrt{n}$ & 3.2 & 31.6 & 1000 & 試除法判斷質數 \\
$n$ & 10 & 1000 & $10^6$ & 掃過陣列一次 \\
$n \lg n$ & 33 & 9966 & $2.0 \times 10^7$ & 合併排序、堆積排序 \\
$n^2$ & 100 & $10^6$ & $10^{12}$ & 兩兩比較、插入排序最壞情況 \\
$n^3$ & 1000 & $10^9$ & $10^{18}$ & 三層迴圈、Floyd--Warshall \\
$2^n$ & 1024 & $1.1 \times 10^{301}$ & 約 $10^{301030}$ & 枚舉所有子集合 \\
$n!$ & $3.6 \times 10^6$ & 約 $10^{2568}$ & 約 $10^{5565709}$ & 枚舉所有排列 \\
\bottomrule
\end{tabular}
\end{center}

\section{分析程式片段}
\label{sec:ana-fragments}

考題和作業常常給一段迴圈，問某一行執行幾次，或問時間複雜度。 做法是從最內層往外
寫成和，再化簡。 下面五段迴圈各自示範一種情況，都數 \texttt{c++} 執行的次數。

\begin{code}{c}{五段迴圈}
// (a) 內層上界跟著外層變
for (int i = 0; i < n; i++)
    for (int j = 0; j < i; j++)
        c++;

// (b) 每次加倍
for (long long i = 1; i < n; i *= 2)
    c++;

// (c) 內層的步長跟著外層變
for (int i = 1; i <= n; i++)
    for (int j = i; j <= n; j += i)
        c++;

// (d) 內層跑到根號
for (int i = 1; i <= n; i++)
    for (long long j = 1; j * j <= i; j++)
        c++;

// (e) 外層每次減半，內層跑 i 次
for (int i = n; i > 0; i /= 2)
    for (int j = 0; j < i; j++)
        c++;
\end{code}

\textbf{(a)} 外層的 $i$ 固定時，內層跑 $i$ 次。 總次數
$\sum_{i=0}^{n-1} i = \frac{n(n-1)}{2} = \Theta(n^2)$。 這和第~\ref{sec:ana-cost}~節的兩
兩比較是同一個和。

\textbf{(b)} $i$ 依序是 $1, 2, 4, \ldots$，只要 $2^k < n$ 就會再跑一次。 執行次數是
滿足 $2^k < n$ 的 $k \ge 0$ 的個數，$n \ge 2$ 時等於 $\lceil \lg n \rceil$，是
$\Theta(\lg n)$。 每次乘以任何大於 $1$ 的常數（例如乘 $3$）都一樣是 $\Theta(\lg n)$，
只差對數的底。

\textbf{(c)} 外層是 $i$ 時，$j$ 是 $i, 2i, 3i, \ldots$，不超過 $n$ 的有 $\lfloor n/i
\rfloor$ 個。 總次數
\[
\sum_{i=1}^{n} \left\lfloor \frac{n}{i} \right\rfloor.
\]
因為 $\frac{n}{i} - 1 < \lfloor \frac{n}{i} \rfloor \le \frac{n}{i}$，總和夾在
$n H_n - n$ 和 $n H_n$ 之間，而 $H_n = \Theta(\lg n)$，所以是 $\Theta(n \lg n)$。 兩
層迴圈不一定是 $n^2$，要看內層實際跑幾次。

\textbf{(d)} 外層是 $i$ 時內層跑 $\lfloor \sqrt{i} \rfloor$ 次。 上界：每一項都不超過
$\sqrt{n}$，總和不超過 $n\sqrt{n}$。 下界：只看 $i \ge n/2$ 的那一半，每一項至少
$\lfloor \sqrt{n/2} \rfloor$，大約有 $n/2$ 項，總和至少是 $n^{3/2}$ 的常數倍。 所以
是 $\Theta(n^{3/2})$。 精確一點的估計是 $\frac{2}{3} n^{3/2}$，來自積分
$\int_0^n \sqrt{x}\,dx$。

「只看後半」是求下界的常用手法：前半的項丟掉，後半的每一項都用最小的那一項代替。
上一節估 $\lg(n!)$ 也是這樣做的。

\textbf{(e)} 外層的 $i$ 依序是 $n, \lfloor n/2 \rfloor, \lfloor n/4 \rfloor, \ldots,
1$，內層跑 $i$ 次，總次數不超過
\[
n + \frac{n}{2} + \frac{n}{4} + \cdots < 2n.
\]
又至少是 $n$（第一輪），所以是 $\Theta(n)$。 這是公比 $1/2$ 的等比級數，和第一項同
一級。 外層雖然跑了 $\Theta(\lg n)$ 輪，但不能直接寫成 $n \lg n$：每一輪的工作量在
減半。

\subsection{實際數一次}

把五段迴圈各包成一個函式，對幾個 $n$ 實際數出次數。 推出公式之後跑一次這種短程式，
很容易抓到推導裡的錯。

\begin{example}[label=ex:loop-counts]{c}{(loop\_counts.c) 五段迴圈的執行次數}
#include <stdio.h>

// 五段迴圈，各自回傳最內層執行的次數。
long long frag_a(int n)
{
    long long c = 0;
    for (int i = 0; i < n; i++)
        for (int j = 0; j < i; j++)
            c++;
    return c;
}

long long frag_b(int n)
{
    long long c = 0;
    for (long long i = 1; i < n; i *= 2)
        c++;
    return c;
}

long long frag_c(int n)
{
    long long c = 0;
    for (int i = 1; i <= n; i++)
        for (int j = i; j <= n; j += i)
            c++;
    return c;
}

long long frag_d(int n)
{
    long long c = 0;
    for (int i = 1; i <= n; i++)
        for (long long j = 1; j * j <= i; j++)
            c++;
    return c;
}

long long frag_e(int n)
{
    long long c = 0;
    for (int i = n; i > 0; i /= 2)
        for (int j = 0; j < i; j++)
            c++;
    return c;
}

int main(void)
{
    int ns[] = {16, 100, 1000, 10000};
    printf("n a b c d e\n");
    for (int k = 0; k < 4; k++) {
        int n = ns[k];
        printf("%d %lld %lld %lld %lld %lld\n", n, frag_a(n),
               frag_b(n), frag_c(n), frag_d(n), frag_e(n));
    }
    return 0;
}
\end{example}
\iosample{}{n a b c d e \\ 16 120 4 50 38 31 \\ 100 4950 7 482 625 197 \\ 1000 499500 10 7069 20615 1994 \\ 10000 49995000 14 93668 661750 19995}

和公式對照（$n = 10^4$）：(a) $n(n-1)/2 = 49{,}995{,}000$，完全一樣；(b)
$\lceil \lg 10^4 \rceil = 14$；(c) 落在 $nH_n - n \approx 87{,}876$ 和
$nH_n \approx 97{,}876$ 之間；(d) 接近 $\frac{2}{3} n^{3/2} \approx 666{,}667$；(e) 小於
$2n = 20{,}000$。 $n$ 從 $10^3$ 變成 $10^4$ 時，(a) 變成 $100$ 倍，(c) 大約 $13$
倍，(d) 大約 $32$ 倍（$10^{1.5} \approx 31.6$），(e) 大約 $10$ 倍，各自符合
$n^2$、$n \lg n$、$n^{1.5}$、$n$ 的成長方式。

\begin{examnote}
問「某一行執行幾次」時，先寫出和（例如 $\sum_{i=1}^{n} \lfloor n/i \rfloor$），再化
成封閉形式或給出上下界，最後才寫 $\Theta$。 寫完後代一個小的 $n$（例如 $n = 4$）
手數一次，和公式比對，可以抓到大部分差 $1$ 的錯誤，例如 $<$ 和 $\le$ 搞混。
\end{examnote}

\section{最好、最壞和平均的情況}
\label{sec:ana-cases}

插入排序的例子說明，同樣大小的輸入，步數可以差很多。 所以談執行時間時，要說清楚
是哪一種：

\begin{itemize}
\item \textbf{最壞情況}（worst case）：所有大小為 $n$ 的輸入裡，步數的最大值。 它
是一個保證：不管輸入是什麼，都不會比這個慢。 這本書沒有特別說明時，執行時間指的
都是最壞情況。
\item \textbf{最好情況}（best case）：步數的最小值。 用處不大，因為輸入通常不會剛
好是最好的那一種，但常出現在選擇題。
\item \textbf{平均情況}（average case）：先假設輸入怎麼分布，再算步數的期望值。
排序常用的假設是：$n$ 個不同的數，$n!$ 種排列出現的機會都一樣。
\end{itemize}

\subsection{插入排序的平均情況}

在「每種排列機會一樣」的假設下，插入排序平均挪動幾次？ 挪動次數等於逆序對的個數，
所以要算逆序對個數的期望值。

任取兩個位置 $i < j$。 把一個排列的第 $i$、$j$ 個數對調，得到另一個排列；這樣兩兩
配成一對，每一對裡剛好一個在 $(i, j)$ 是逆序。 所以 $(i, j)$ 是逆序的排列剛好佔一
半，機率是 $1/2$。 一共有 $\binom{n}{2}$ 對位置，期望值可以一對一對加起來（期望值
的線性），得到
\[
\text{平均挪動次數} = \binom{n}{2} \cdot \frac{1}{2} = \frac{n(n-1)}{4}.
\]
這剛好是最壞情況的一半，仍然是 $\Theta(n^2)$。

$n$ 很小時，可以把 $n!$ 種排列全部列出來，直接驗證這個公式。 下面的程式用遞迴產
生 $0, 1, \ldots, n-1$ 的所有排列（《C 程式設計基礎》第~15~章回溯法的寫法），對每
一個排列跑一次插入排序，加總挪動次數。

\begin{example}[label=ex:avg-shift]{c}{(avg\_shift.c) 所有排列的平均挪動次數}
#include <stdio.h>

#define MAXN 8

int perm[MAXN];
int used[MAXN];
long long total;        // 所有排列的挪動次數總和
long long count;        // 排列個數

// 插入排序挪動的次數，等於 a[0..n-1] 裡 i < j 且 a[i] > a[j] 的對數。
// 這裡直接把插入排序跑一次來數，不用這個性質。
long long shifts(const int *src, int n)
{
    int a[MAXN];
    long long s = 0;
    for (int i = 0; i < n; i++)
        a[i] = src[i];
    for (int i = 1; i < n; i++) {
        int key = a[i];
        int j = i - 1;
        while (j >= 0 && a[j] > key) {
            a[j + 1] = a[j];
            s++;
            j--;
        }
        a[j + 1] = key;
    }
    return s;
}

// 依序填 perm[pos..n-1]，填滿時算一次。
void gen(int pos, int n)
{
    if (pos == n) {
        total += shifts(perm, n);
        count++;
        return;
    }
    for (int v = 0; v < n; v++) {
        if (!used[v]) {
            used[v] = 1;
            perm[pos] = v;
            gen(pos + 1, n);
            used[v] = 0;
        }
    }
}

int main(void)
{
    printf("n n! total average n(n-1)/4\n");
    for (int n = 1; n <= MAXN; n++) {
        total = 0;
        count = 0;
        gen(0, n);
        printf("%d %lld %lld %.2f %.2f\n", n, count, total,
               (double)total / count, n * (n - 1) / 4.0);
    }
    return 0;
}
\end{example}
\iosample{}{n n! total average n(n-1)/4 \\ 1 1 0 0.00 0.00 \\ 2 2 1 0.50 0.50 \\ 3 6 9 1.50 1.50 \\ 4 24 72 3.00 3.00 \\ 5 120 600 5.00 5.00 \\ 6 720 5400 7.50 7.50 \\ 7 5040 52920 10.50 10.50 \\ 8 40320 564480 14.00 14.00}

八個 $n$ 都和公式一致。 這種「小範圍全部列舉」的程式證明不了公式，但公式推錯時
通常馬上會露出來。

\subsection{「插入排序是 $\Theta(n^2)$」對不對}

這句話常在考卷的是非題出現，答案要看它指的是什麼。

\begin{itemize}
\item 「插入排序最壞情況的執行時間是 $\Theta(n^2)$」：對。
\item 「插入排序的執行時間是 $O(n^2)$」：對。 每一種輸入都不超過最壞情況，所以
$O(n^2)$ 是對所有輸入都成立的上界。
\item 「插入排序的執行時間是 $\Omega(n)$」：對。 每一種輸入至少要把每個數看一次。
\item 「插入排序的執行時間是 $\Theta(n^2)$」，沒有說哪一種情況：不對。 排好的輸入
只要 $\Theta(n)$，所以不是每一種輸入都至少 $n^2$ 的常數倍。
\item 「插入排序的執行時間是 $\Omega(n^2)$」：同上，不對。
\end{itemize}

整理成一句：對所有輸入成立的上界，看最壞情況；對所有輸入成立的下界，看最好情況。
「演算法 X 是 $\Theta(f(n))$」只有在最好和最壞情況同一級時才能不加說明。

\section{輸入大小要用位元數來算的時候}
\label{sec:ana-input-size}

前面的例子輸入是 $n$ 個數，大小就是 $n$。 如果輸入只有一個整數 $x$，問 $x$ 是不是
質數呢？ 試除法從 $2$ 開始試，試到 $\sqrt{x}$ 為止，步數大約是 $\sqrt{x}$。 這看起
來比線性還快，但輸入的大小並不是 $x$。 把 $x$ 寫下來只需要 $b = \lfloor \lg x
\rfloor + 1$ 個位元（或 $d$ 個十進位位數），輸入大小應該是 $b$。 用 $b$ 來寫，
步數是 $\sqrt{x} \approx 2^{b/2}$，是 $b$ 的指數函數。

\begin{example}[label=ex:trial-div]{c}{(trial\_div.c) 試除法的步數和位數}
#include <stdio.h>

// 試除法判斷質數。*steps 記錄做了幾次除法（取餘數）。
int is_prime(long long x, long long *steps)
{
    *steps = 0;
    if (x < 2) {
        return 0;
    }
    for (long long d = 2; d * d <= x; d++) {
        (*steps)++;
        if (x % d == 0) {
            return 0;
        }
    }
    return 1;
}

int main(void)
{
    long long xs[] = {101, 10007, 1000003, 100000007, 10000000019LL,
                      1000000000039LL};
    printf("x digits prime steps\n");
    for (int k = 0; k < 6; k++) {
        long long steps;
        int p = is_prime(xs[k], &steps);
        int digits = 0;
        for (long long t = xs[k]; t > 0; t /= 10)
            digits++;
        printf("%lld %d %s %lld\n", xs[k], digits,
               p ? "yes" : "no", steps);
    }
    return 0;
}
\end{example}
\iosample{}{x digits prime steps \\ 101 3 yes 9 \\ 10007 5 yes 99 \\ 1000003 7 yes 999 \\ 100000007 9 yes 9999 \\ 10000000019 11 yes 99999 \\ 1000000000039 13 yes 999999}

六個數都是質數，所以每次都要試到最後。 位數每多 $2$，步數就變成 $10$ 倍：輸入只
長了一點，工作量卻乘上一個固定倍數，這是指數成長的樣子。 $30$ 位數的質數，試除法
要試大約 $10^{15}$ 次。

像這樣執行時間是輸入\textbf{數值}的多項式、卻是輸入\textbf{位元數}的指數的演算
法，叫做\textbf{擬多項式時間}（pseudo-polynomial time）演算法。 第~18~章的 0/1 背
包動態規劃要 $\Theta(nW)$ 的時間，$W$ 是背包容量的數值，也是擬多項式；第~28~章談
NP 完備時，會說明為什麼這個區別很重要。

RAM 模型假設一個字組放得下 $n$ 這個量級的數，所以排序 $n$ 個 \texttt{int} 時，比
較兩個數算一步是合理的。 但如果數字本身有幾千位（例如密碼學裡的大整數），加法和
乘法就不能算一步，要另外數位元運算的次數。

\section{本章常見錯誤}
\label{sec:ana-pitfalls}

\begin{center}
\begin{tabular}{>{\raggedright\arraybackslash}p{0.24\linewidth}>{\raggedright\arraybackslash}p{0.50\linewidth}p{0.16\linewidth}}
\toprule
錯誤 & 例子 & 見 \\
\midrule
把 $O$ 當成「剛好是」 & 以為 $n = O(n^2)$ 是錯的；以為「$O(n^2)$ 的演算法一定比
$O(n)$ 的慢」 & 第~\ref{sec:ana-notation}~節 \\
沒說是哪一種情況 & 「插入排序是 $\Theta(n^2)$」不加說明 &
第~\ref{sec:ana-cases}~節 \\
證明只代幾個 $n$ & 用定義證明 $\Theta$ 時沒有寫出 $c_1$、$c_2$、$n_0$，或只驗證
$n = 1, 2, 3$ & 第~\ref{sec:ana-notation}~節 \\
兩層迴圈就寫 $n^2$ & 內層步長是 $i$ 的迴圈（$n \lg n$），外層減半的迴圈（$n$）&
第~\ref{sec:ana-fragments}~節 \\
指數的底當成常數丟掉 & 以為 $4^n = O(2^n)$；以為 $2^{2\lg n} = \Theta(n)$ &
第~\ref{sec:ana-functions}~節 \\
極限不存在就下結論 & 用 $\lim f/g$ 判斷，卻沒檢查極限存在 &
第~\ref{sec:ana-notation}~節 \\
輸入大小量錯 & 說試除法是 $O(\sqrt{x})$ 所以是「多項式時間」 &
第~\ref{sec:ana-input-size}~節 \\
C 和虛擬碼的編號差 $1$ & 虛擬碼的 $j \ge 1$ 翻成 \texttt{j >= 1}，少處理了
\texttt{a[0]} & 第~\ref{sec:ana-conventions}~節 \\
\bottomrule
\end{tabular}
\end{center}

\section{觀念題}
\label{sec:ana-exercises}

題型和研究所考試、期中期末考的筆試相近，有選擇、計算、追蹤和證明。 解答在本章最後。

\begin{enumerate}
\item （選擇，可能不只一個答案）下列哪些敘述正確？
\begin{enumerate}
\item $100 n \lg n = O(n^2)$
\item $n^2 = O(100 n \lg n)$
\item $2^{n+3} = \Theta(2^n)$
\item $3^n = O(2^n)$
\item $\lg(n^5) = \Theta(\lg n)$
\end{enumerate}

\item （計算）求下面的程式片段中 \texttt{x++} 執行的確切次數，以 $n$ 表示，並寫出
$\Theta$。
\begin{code}{c}{觀念題 2}
for (int i = 1; i <= n; i++)
    for (int j = 1; j <= i; j++)
        for (int k = j; k <= i; k++)
            x++;
\end{code}

\item （計算）求下面的程式片段的時間複雜度，以 $\Theta$ 表示，並說明理由。
\begin{code}{c}{觀念題 3}
for (int i = n; i >= 1; i /= 3)
    for (long long j = 1; j <= (long long)i * i; j++)
        x++;
\end{code}

\item （計算）$n$ 是 $2$ 的次方。 求 \texttt{x++} 執行的確切次數。
\begin{code}{c}{觀念題 4}
for (int i = 1; i <= n; i *= 2)
    for (int j = 1; j <= i; j++)
        x++;
\end{code}

\item （排序）把下列函數依成長速度由慢到快排列。 同一級（彼此是 $\Theta$）的用
等號連起來。
\[
n^{1.5},\quad n^{\lg 3},\quad (\lg n)^{\lg n},\quad 2^{\lg n},\quad \lg(n!),\quad
4^{\lg n},\quad n^{1/\lg n},\quad \sqrt{n}\,\lg n,\quad n \lg n,\quad (1.5)^n
\]

\item （證明）用定義證明 $5n^2 - 12n + 4 = \Theta(n^2)$，寫出 $c_1$、$c_2$、$n_0$。

\item （證明或反例）$f$、$g$、$h$ 都是正的函數。 下列敘述對的請證明，錯的請給反例。
\begin{enumerate}
\item 若 $f(n) = O(g(n))$，則 $f(n)^2 = O(g(n)^2)$。
\item 若 $f(n) = O(h(n))$ 且 $g(n) = O(h(n))$，則 $f(n) \cdot g(n) = O(h(n)^2)$。
\item 若 $f(n) = O(g(n))$，則 $\lim_{n \to \infty} f(n)/g(n)$ 存在。
\item 若 $f(n) = O(g(n))$，則 $f(n) - g(n) = O(1)$。
\end{enumerate}

\item （證明）用定義~\ref{def:little-o} 或定理~\ref{thm:limit} 證明
$n \lg n = o(n^{1.1})$，並說明為什麼 $n^{1.1} \ne O(n \lg n)$。

\item \label{ana-q-trace}（追蹤）用插入排序（虛擬碼~\ref{alg:insertion}）排序
$\langle 50, 20, 40, 10, 70, 30, 60 \rangle$。 寫出外層迴圈每一輪之後的陣列，以及
$A[j] > \mathit{key}$ 的總比較次數和總挪動次數。

\item （計算）插入排序排一個長度 $n$、有 $k$ 個逆序對的陣列。 證明比較次數
（只數 $A[j] > \mathit{key}$）等於 $k + (n - 1) - z$，其中 $z$ 是「$\mathit{key}$
被放到第 $1$ 格」的輪數。 用第~\ref{ana-q-trace}~題的結果驗證這個公式。

\item \label{ana-q-lower-bound}（證明）下面的程序在已經由小到大排好的 $A[1..n]$ 裡，回傳第一個大於等於
$x$ 的位置；全部都小於 $x$ 時回傳 $n + 1$。 寫出 \textbf{while} 迴圈的不變量，用
它證明程序正確，並說明迴圈一定會停。
\begin{pseudo}[unbreakable]{第一個大於等於 $x$ 的位置}
\Function{Lower-Bound}{$A, n, x$}
    \State $\mathit{lo} \gets 1$
    \State $\mathit{hi} \gets n + 1$
    \While{$\mathit{lo} < \mathit{hi}$}
        \State $\mathit{mid} \gets \lfloor (\mathit{lo} + \mathit{hi}) / 2 \rfloor$
        \If{$A[\mathit{mid}] < x$}
            \State $\mathit{lo} \gets \mathit{mid} + 1$
        \Else
            \State $\mathit{hi} \gets \mathit{mid}$
        \EndIf
    \EndWhile
    \State \Return $\mathit{lo}$
\EndFunction
\end{pseudo}

\item （選擇）在 $n$ 個互不相同的數裡從左到右找 $x$，找到就停。 已知 $x$ 一定在
裡面，而且在每個位置的機率都是 $1/n$。 平均比較幾次？
(A) $n/2$ \quad (B) $(n+1)/2$ \quad (C) $n$ \quad (D) $\lg n$ \quad (E) $(n-1)/2$

\item （是非）下面的函式用迴圈算 $1 + 2 + \cdots + x$，迴圈跑 $x$ 次。 「因為執行
時間是 $\Theta(x)$，所以它是多項式時間的演算法。」 這句話對嗎？ 說明理由。
\begin{code}{c}{觀念題 13}
long long sum_to(long long x)
{
    long long s = 0;
    for (long long i = 1; i <= x; i++)
        s += i;
    return s;
}
\end{code}

\item （證明）證明 $f(n) + g(n) = \Theta(\max(f(n), g(n)))$ 不一定能推廣到
「$n$ 個函數」：給一個例子，$f_1, f_2, \ldots$ 都是 $\Theta(1)$，但
$\sum_{i=1}^{n} f_i(n)$ 不是 $\Theta(1)$。 再說明這和第~\ref{sec:ana-fragments}~節
(e) 不能寫成「$\Theta(\lg n)$ 輪，每輪 $O(n)$，所以 $\Theta(n \lg n)$」有什麼關係。
\end{enumerate}

\section{歷屆試題}
\label{sec:ana-pastexam}

以下是研究所入學考試的原題，題目照原卷（語言、記號、配分照錄）。 詳解在本章最後。

\begin{pastexam}[label={px:ana-1}]{臺大資工所 115 學年度 第 2 題}
(4 points) For positive functions $f(n)$ and $g(n)$ on $\mathbb{N}$. If $f(n)=O(g(n))$
and $f(n)=\Omega(1)$, how many of the following statements are true?
\begin{itemize}
\item $g(n)=\Omega(f(n))$
\item $f(n)\cdot\log(1+f(n))=O(g(n)\cdot\log(1+g(n)))$
\item $2^{f(n)}=O(2^{g(n)})$
\item $(f(n))^k=O((g(n))^k)$ for any $0<k<1$
\end{itemize}
(A) 0 \qquad (B) 1 \qquad (C) 2 \qquad (D) 3 \qquad (E) 4
\end{pastexam}

\begin{pastexam}[label={px:ana-2}]{臺科大資工所 113 學年度 第 I-1 題}
I.\quad (35\%) Multiple choice questions.

(5\%) Order the following functions by growth rate in increasing order.\\
$N$, $\sqrt{N}$, $N^{1.5}$, $N\log N$, $N^2$, $N\log(\log N)$, $N\log^2 N$,
$N\log(N^2)$, $2^N$, 2048\\
Which is the 5th function in the ordered sequence?\\
(a) $N$,\quad (b) $\sqrt{N}$,\quad (c) $N^{1.5}$,\quad (d) $N\log N$,\quad (e) $N^2$,\\
(f) $N\log(\log N)$,\quad (g) $N\log^2 N$,\quad (h) $N\log(N^2)$,\quad (i) $2^N$,\quad
(j) 2048
\end{pastexam}

\begin{pastexam}[label={px:ana-3}]{臺科大電子所甲組 110 學年度 第 4 題}
Suppose we have three functions:
\begin{minted}[fontsize=\small]{text}
foo(int N)
    do something in A steps
    for i equals 1 to N
        do something in B steps
        for j equals 1 to N
            do something in C steps

goo(int N)
    do something in A steps
    for i equals 1 to N
        do something in B step
        for j equals 1 to i
            do something in C steps

poo(int N)
    do something in A steps
    for i equals 1 to N
        do something in B steps
        for j equals 1 to 250
            do something in C steps
\end{minted}
\begin{enumerate}[label=(\alph*)]
\item (6\%) For each of the three methods above (foo(), goo(), and poo()), what is the
exact runtime for each method (in steps) in terms of A, B, C, and N? You should not count
loop related operations.
\item (4\%) For N = 100, which one is the fastest algorithm? Why?
\item (6\%) Assuming that A, B, and C are constants, what is the asymptotic runtime of
each method in terms of N?
\item (4\%) As N grows infinitely large, which one is the fastest algorithm?
\end{enumerate}
\end{pastexam}

\begin{pastexam}[label={px:ana-4}]{中山資工所 112 學年度 第 2 題}
The pseudo code for the insertion sort is given as follows.
\begin{minted}[fontsize=\small]{text}
Read input elements into a[ ]; //stored in a[1], a[2], ..., a[n]
for i = 2 to n do
{
   j = i;
   while (j >= 2) and  (a[j] < a[j-1]) do
   {
      Swap(a[j], a[j-1]); // exchange a[j] and a[j-1]
      j = j-1;
   }
}
\end{minted}
\begin{enumerate}[label=(\alph*)]
\item (5\%) Suppose that the input elements are 3, 9, 6, 5, 8, 2. How many times
(counts) are needed for the execution of \texttt{Swap( )}?
\item (5\%) Suppose that the input is a permutation of $\{1,2,3,\ldots,n\}$. What
permutation will make \texttt{Swap( )} be executed the least times? How many times?
\item (5\%) Suppose that the input is a permutation of $\{1,2,3,\ldots,n\}$. What
permutation will make \texttt{Swap( )} be executed the most times? How many times?
\end{enumerate}
\end{pastexam}

\begin{pastexam}[label={px:ana-5}]{政大資科所 110 學年度 第 1 題}
(10\%) Prove or disprove the following statements ($n$ is the size of the input array):
\begin{enumerate}[label=\alph*.]
\item (5\%) The time complexity of any sorting algorithm is $O(n^2)$.
\item (5\%) The best-case time complexity of insertion sort is $\theta(n\log n)$.
\end{enumerate}
\end{pastexam}

\section{實作練習}
\label{sec:ana-practice}

下面是 Aki-Judge 上和本章相關的題目。 這一章的重點是分析，能寫成程式題的不多，
其他練習請以觀念題為主。

\begin{center}
\begin{tabular}{llc}
\toprule
小節 & 題目 & 難度 \\
\midrule
分析程式片段 & \akiref{count-nested-loop-executions}{計算巢狀迴圈執行次數} & 暖身 \\
插入排序 & \akiref{insertion-sort-function}{插入排序函式} & 基本 \\
\bottomrule
\end{tabular}
\end{center}

「插入排序函式」只要交排序函式，不寫 \texttt{main}。 交之前先用
範例~\ref{ex:insertion-count} 的三組輸入自己測一次，特別是已經排好和完全相反的兩組。

\section{讀完你應該會}
\label{sec:ana-checklist}

\begin{itemize}
\item 看到一段虛擬碼或 C 程式，能逐行寫出執行次數，化成封閉形式或上下界，再寫成
$\Theta$。
\item 能用定義證明 $f(n) = \Theta(g(n))$，寫出 $c_1$、$c_2$、$n_0$；能用反證法說明
$f(n) \ne O(g(n))$；能用極限判斷，並知道極限不存在時要回到定義。
\item 能把對數、多項式、指數、階乘混在一起的一串函數依成長速度排好，包括
$a^{\log_b n} = n^{\log_b a}$ 和 $\lg(n!) = \Theta(n \lg n)$ 這類換算。
\item 能用迴圈不變量的三步（初始、維持、結束）證明一個迴圈的正確性。
\item 能分別說出插入排序最好、最壞、平均情況的次數，並判斷「某演算法是
$\Theta(f(n))$」這類敘述在哪一種意義下成立。
\item 輸入是一個整數時，能用位元數判斷一個演算法是不是多項式時間。
\end{itemize}

\textbf{延伸閱讀}（書末會整理完整的參考資料）

\begin{itemize}
\item Thomas H. Cormen 等，\textit{Introduction to Algorithms}，第 4 版，第 2、3
章與附錄 A：插入排序的逐行分析、迴圈不變量、漸近符號的完整定義；CLRS 的附錄 A 整理了求
和的技巧（本書的求和工具見附錄~A）。
\item Steven S. Skiena，\textit{The Algorithm Design Manual}，第 3 版，第 2 章：
用較少的符號講 RAM 模型和成長速度，例子偏向實務。
\item Antti Laaksonen，\textit{Competitive Programmer's Handbook}，第 2 章：幾頁
就講完常見的複雜度和「資料量對應可接受的複雜度」，適合快速複習。
\end{itemize}

\addcontentsline{toc}{section}{觀念題解答}
\section*{觀念題解答}

\begin{enumerate}
\item (a)、(c)、(e) 正確。 (a)：$\lg n \le n$，所以 $100n\lg n \le 100n^2$。
(b)：$n^2 / (100 n \lg n) = n / (100 \lg n) \to \infty$，不是 $O$。 (c)：
$2^{n+3} = 8 \cdot 2^n$，常數倍。 (d)：$3^n / 2^n = 1.5^n \to \infty$。 (e)：
$\lg(n^5) = 5\lg n$。

\item 固定 $i$、$j$ 時，最內層跑 $i - j + 1$ 次。 固定 $i$ 時，
$\sum_{j=1}^{i} (i - j + 1) = 1 + 2 + \cdots + i = \frac{i(i+1)}{2}$。 總次數
\[
\sum_{i=1}^{n} \frac{i(i+1)}{2}
= \frac{1}{2}\left(\frac{n(n+1)(2n+1)}{6} + \frac{n(n+1)}{2}\right)
= \frac{n(n+1)(n+2)}{6}.
\]
是 $\Theta(n^3)$。 驗算 $n = 2$：$i = 1$ 時 $1$ 次，$i = 2$ 時 $j = 1$ 有 $2$ 次、
$j = 2$ 有 $1$ 次，共 $4 = \frac{2 \cdot 3 \cdot 4}{6}$。

\item 外層的 $i$ 是 $n, \lfloor n/3 \rfloor, \lfloor n/9 \rfloor, \ldots$，內層跑
$i^2$ 次。 總和不超過 $n^2 + n^2/9 + n^2/81 + \cdots = n^2 \cdot \frac{1}{1 - 1/9} =
\frac{9}{8}n^2$，又至少是第一輪的 $n^2$，所以是 $\Theta(n^2)$。 公比 $1/9 < 1$，總
和和第一項同一級。

\item 設 $n = 2^m$。 $i = 1, 2, 4, \ldots, 2^m$，內層跑 $i$ 次，總次數
$1 + 2 + \cdots + 2^m = 2^{m+1} - 1 = 2n - 1$，是 $\Theta(n)$。

\item 先化簡：$2^{\lg n} = n$，$4^{\lg n} = n^2$，$n^{1/\lg n} = 2^{\lg n / \lg n} = 2$，
$(\lg n)^{\lg n} = 2^{\lg n \cdot \lg\lg n} = n^{\lg \lg n}$，$\lg(n!) = \Theta(n \lg n)$，
$n^{\lg 3} \approx n^{1.585}$。 由慢到快：
\[
n^{1/\lg n} < \sqrt{n}\,\lg n < 2^{\lg n} < \lg(n!) = n\lg n < n^{1.5} < n^{\lg 3}
< 4^{\lg n} < (\lg n)^{\lg n} < (1.5)^n.
\]
$(\lg n)^{\lg n} = n^{\lg\lg n}$ 的次方 $\lg \lg n$ 會無限變大，所以比任何多項式都快；
和 $(1.5)^n$ 比，取 $\lg$ 後是 $\lg n \cdot \lg \lg n$ 對 $n \lg 1.5$，後者大。

\item 上界：$n \ge 1$ 時 $-12n + 4 \le -8 < 0$，所以 $5n^2 - 12n + 4 \le 5n^2$，取
$c_2 = 5$。 下界：要 $5n^2 - 12n + 4 \ge 2n^2$，即 $3n^2 - 12n + 4 \ge 0$。 這個二次
式的較大根是 $\frac{12 + \sqrt{96}}{6} \approx 3.63$，所以 $n \ge 4$ 時成立（$n = 4$：
$48 - 48 + 4 = 4 \ge 0$）。 取 $c_1 = 2$、$c_2 = 5$、$n_0 = 4$，此時 $2n^2 \ge 0$ 也
成立。

\item
\begin{enumerate}
\item 對。 $n \ge n_0$ 時 $0 \le f(n) \le c\,g(n)$，兩邊平方得
$f(n)^2 \le c^2 g(n)^2$，取常數 $c^2$。
\item 對。 $f \le c_1 h$（$n \ge n_1$），$g \le c_2 h$（$n \ge n_2$），相乘得
$fg \le c_1 c_2 h^2$，$n \ge \max(n_1, n_2)$。
\item 錯。 $g(n) = n$，$f(n)$ 在偶數時是 $n$、奇數時是 $1$。 $f(n) \le g(n)$，所以
$f = O(g)$；但 $f/g$ 在 $1$ 和 $1/n$ 之間跳，沒有極限。
\item 錯。 $f(n) = 2n$，$g(n) = n$。 $f = O(g)$（取 $c = 2$），但 $f(n) - g(n) = n$
不是 $O(1)$。 $O$ 允許差一個常數倍，不是差一個常數。
\end{enumerate}

\item $\dfrac{n \lg n}{n^{1.1}} = \dfrac{\lg n}{n^{0.1}}$。 用羅必達法則（把 $\lg n$
寫成 $\ln n / \ln 2$）：$\dfrac{1/(n \ln 2)}{0.1\, n^{-0.9}} = \dfrac{10}{\ln 2 \cdot
n^{0.1}} \to 0$。 由定理~\ref{thm:limit}，$n \lg n = o(n^{1.1})$。 同一個定理的第 1 點
也說明 $n \lg n \ne \Omega(n^{1.1})$，由轉置性質，$n^{1.1} \ne O(n \lg n)$。 注意
$n = 2^{60}$ 時 $\lg n = 60$，$n^{0.1} = 64$，兩者才差不多大；$n^{1.1}$ 要在非常
大的 $n$ 才明顯超過 $n \lg n$。

\item 每一輪之後的陣列：
\begin{center}
\begin{tabular}{cccl}
\toprule
$i$ & 比較 & 挪動 & 這一輪之後 \\
\midrule
2 & 1 & 1 & 20 50 40 10 70 30 60 \\
3 & 2 & 1 & 20 40 50 10 70 30 60 \\
4 & 3 & 3 & 10 20 40 50 70 30 60 \\
5 & 1 & 0 & 10 20 40 50 70 30 60 \\
6 & 4 & 3 & 10 20 30 40 50 70 60 \\
7 & 2 & 1 & 10 20 30 40 50 60 70 \\
\midrule
合計 & 13 & 9 & \\
\bottomrule
\end{tabular}
\end{center}
驗證的方法：把範例~\ref{ex:insertion-count} 的輸入改成三行
\texttt{1}、\texttt{7}、\texttt{50 20 40 10 70 30 60}，程式印出
\texttt{comparisons = 13, shifts = 9}。

\item 第 $i$ 輪挪動 $s_i$ 次。 如果這一輪因為某個 $A[j] \le \mathit{key}$ 而停下，
比較了 $s_i + 1$ 次；如果因為 $j = 0$ 停下（$\mathit{key}$ 放到第 $1$ 格），最後那
次沒有比較，比較了 $s_i$ 次。 把 $n - 1$ 輪加起來：比較次數 $= \sum s_i + (n - 1) - z
= k + (n - 1) - z$，因為挪動總次數等於逆序對個數 $k$。 第~\ref{ana-q-trace}~題
$k = 9$、$n = 7$，
$\mathit{key}$ 被放到第 $1$ 格的是 $i = 2$（20）和 $i = 4$（10），$z = 2$，
$9 + 6 - 2 = 13$，和表上一樣。

\item 不變量：每次檢查 \textbf{while} 條件時，$1 \le \mathit{lo} \le \mathit{hi} \le
n + 1$，而且 $A[1..\mathit{lo}-1]$ 的每個數都小於 $x$，$A[\mathit{hi}..n]$ 的每個數都
大於等於 $x$。

初始：$\mathit{lo} = 1$、$\mathit{hi} = n + 1$，兩段都是空的，成立。

維持：$\mathit{lo} < \mathit{hi}$ 時 $\mathit{lo} \le \mathit{mid} < \mathit{hi}$，所以
$A[\mathit{mid}]$ 存在。 若 $A[\mathit{mid}] < x$，因為陣列排好了，
$A[\mathit{lo}..\mathit{mid}]$ 都小於 $x$，設 $\mathit{lo} \gets \mathit{mid} + 1$ 後第
一段仍然都小於 $x$，而且 $\mathit{lo} \le \mathit{hi}$。 否則 $A[\mathit{mid}] \ge x$，
$A[\mathit{mid}..n]$ 都大於等於 $x$，設 $\mathit{hi} \gets \mathit{mid}$ 後第二段仍然
成立，而且 $\mathit{lo} \le \mathit{mid} = \mathit{hi}$。

結束：迴圈停在 $\mathit{lo} = \mathit{hi}$。 這時 $A[1..\mathit{lo}-1]$ 都小於 $x$、
$A[\mathit{lo}..n]$ 都大於等於 $x$，所以 $\mathit{lo}$ 正是第一個大於等於 $x$ 的位置
（$\mathit{lo} = n + 1$ 表示全部都小於 $x$）。

會停：每一輪 $\mathit{hi} - \mathit{lo}$ 都嚴格變小（$\mathit{lo}$ 變成
$\mathit{mid} + 1 > \mathit{lo}$，或 $\mathit{hi}$ 變成 $\mathit{mid} < \mathit{hi}$），
而它是非負整數，所以有限輪後一定變成 $0$。 每一輪大約減半，共 $\Theta(\lg n)$ 輪。

\item (B)。 $x$ 在第 $k$ 個位置時比較 $k$ 次，平均
$\frac{1}{n}\sum_{k=1}^{n} k = \frac{n+1}{2}$。

\item 不對。 輸入是一個整數 $x$，大小是它的位元數 $b \approx \lg x$。 執行時間
$\Theta(x) = \Theta(2^b)$，是輸入大小的指數，這是擬多項式時間。 （用公式
$x(x+1)/2$ 可以在 $O(1)$ 次算術運算內算完，$x$ 很大時要算大整數的乘法，時間是
$b$ 的多項式。）

\item 取 $f_i(n) = 1$。 每一個都是 $\Theta(1)$，但 $\sum_{i=1}^{n} f_i(n) = n$ 不是
$\Theta(1)$。 「加法」性質只對固定個數的函數成立；項數跟著 $n$ 變多時，要真的把和
算出來。 (e) 的情況相反：輪數 $\Theta(\lg n)$，每一輪的上界 $O(n)$ 都對，所以
$O(n \lg n)$ 是對的上界，但每一輪的實際工作量在減半，把和算出來只有 $\Theta(n)$。
「輪數乘上每輪的上界」只能得到上界，要得到 $\Theta$ 必須算和。
\end{enumerate}

\addcontentsline{toc}{section}{歷屆試題詳解}
\section*{歷屆試題詳解}

\pastsol{px:ana-1}{臺大資工所 115 學年度 第 2 題}
答案是 (D)：第 1、2、4 句對，第 3 句錯。

先把已知條件照定義~\ref{def:big-o} 寫出來。 存在 $c > 0$、$n_1$，使 $n \ge n_1$ 時
$f(n) \le c\,g(n)$；存在 $\delta > 0$、$n_2$，使 $n \ge n_2$ 時 $f(n) \ge \delta$。 把
$c$ 調大不會讓不等式變錯，所以不妨設 $c \ge 1$。 以下都假設 $n \ge \max(n_1, n_2)$，
$\log$ 的底是任何大於 $1$ 的固定數。

\begin{itemize}
\item 第 1 句對。 這是第~\ref{sec:ana-notation}~節的轉置性質：$f(n) \le c\,g(n)$ 就是
$\frac{1}{c} f(n) \le g(n)$，取常數 $1/c$ 得到 $g(n) = \Omega(f(n))$。
\item 第 2 句對。 先估 $\log(1 + f)$。 因為 $c \ge 1$，
$1 + f \le 1 + c\,g \le c\,(1 + g)$，所以 $\log(1 + f) \le \log c + \log(1 + g)$。 要
把常數 $\log c$ 併進 $\log(1 + g)$，需要 $\log(1 + g)$ 不會趨近 $0$。 這裡用到
$f = \Omega(1)$：$g \ge f / c \ge \delta / c$，所以
$\log(1 + g) \ge \log(1 + \delta / c) = \lambda$，$\lambda$ 是一個正的常數。 於是
$\log c \le \frac{\log c}{\lambda} \log(1 + g)$，
\[
\log(1 + f) \le \Bigl(1 + \frac{\log c}{\lambda}\Bigr) \log(1 + g).
\]
再和 $f \le c\,g$ 相乘（兩邊都是正的），得到
$f \log(1 + f) \le c\,\bigl(1 + \frac{\log c}{\lambda}\bigr)\, g \log(1 + g)$，括號外
是常數。
\item 第 3 句錯。 反例：$f(n) = 2n$、$g(n) = n$。 $f = O(g)$（取 $c = 2$），
$f = \Omega(1)$，但 $2^{f(n)} / 2^{g(n)} = 2^{n} \to \infty$，由定理~\ref{thm:limit}，
$2^{2n} \ne O(2^n)$。 第~\ref{sec:ana-functions}~節說過，指數上的常數倍不能丟。
\item 第 4 句對。 $k > 0$ 時 $x^k$ 在正數上是遞增函數，$f \le c\,g$ 兩邊取 $k$ 次方得
$f^k \le c^k g^k$，取常數 $c^k$。 這一步其實對任何 $k > 0$ 都成立。
\end{itemize}

\pastsol{px:ana-2}{臺科大資工所 113 學年度 第 I-1 題}
答案是 (d) $N \log N$。

先化簡：$N \log(N^2) = 2N \log N$，$N \log^2 N = N (\log N)^2$。 $\log$ 的底不影響比
較（第~\ref{sec:ana-functions}~節）。 由小到大是
\[
2048 < \sqrt{N} < N < N \log(\log N) < N \log N < N \log(N^2) < N \log^2 N < N^{1.5}
< N^2 < 2^N .
\]
相鄰兩個逐一說明（除了第 5、6 個，每一對的比值極限都是 $0$，用定理~\ref{thm:limit}）：
\begin{itemize}
\item $2048$ 是常數，$\sqrt{N} \to \infty$；$\sqrt{N} / N = 1/\sqrt{N} \to 0$。
\item $N / (N \log\log N) = 1 / \log\log N \to 0$。
\item $N \log\log N / (N \log N) = \log\log N / \log N$：令 $m = \log N$，這是
$\log m / m \to 0$。
\item $N \log N$ 和 $N \log(N^2) = 2N \log N$ 只差 $2$ 倍，是同一個 $\Theta$；但
$N \ge 2$ 時 $N \log N$ 嚴格比較小。
\item $2N \log N / (N \log^2 N) = 2 / \log N \to 0$。
\item $N \log^2 N / N^{1.5} = (\log N)^2 / N^{0.5} \to 0$（第~\ref{sec:ana-functions}~節
第 1 點：任何次方的對數都比任何正次方的多項式慢）。
\item $N^{1.5} / N^2 \to 0$；$N^2 / 2^N \to 0$（第 2 點）。
\end{itemize}
第 5 個是 $N \log N$。 題目要的是「排好的序列裡的第 5 個」，第 5、6 個同一級時，照
實際大小排，$N \log N$ 在前。 （代 $N = 2^{64}$ 實際算出十個值再排序，順序和上面相
同。）

\pastsol{px:ana-3}{臺科大電子所甲組 110 學年度 第 4 題}
題目說不計迴圈本身的操作，所以只數三種 do something 各做幾次，再乘上各自的步數，
和第~\ref{sec:ana-line-cost}~節的逐行計算一樣。 原題 goo 第 4 行的 ``B step'' 應該是
``B steps''，不影響答案。

(a)
\begin{itemize}
\item foo：A 做 $1$ 次；外層 $i$ 從 $1$ 到 $N$，B 做 $N$ 次；內層每輪跑 $N$ 次，C 做
$N \cdot N$ 次。 $T_{\mathrm{foo}}(N) = A + NB + N^2 C$。
\item goo：內層 $j$ 從 $1$ 到 $i$，跑 $i$ 次，C 共做
$\sum_{i=1}^{N} i = N(N+1)/2$ 次。 $T_{\mathrm{goo}}(N) = A + NB + N(N+1)C/2$。
\item poo：內層固定跑 $250$ 次，C 共做 $250N$ 次。
$T_{\mathrm{poo}}(N) = A + NB + 250NC$。
\end{itemize}

(b) 代 $N = 100$：
\[
T_{\mathrm{foo}} = A + 100B + 10000C,\quad
T_{\mathrm{goo}} = A + 100B + 5050C,\quad
T_{\mathrm{poo}} = A + 100B + 25000C .
\]
三者的 A、B 部分相同，只比 C 的次數：$5050 < 10000 < 25000$，所以 goo 最快（假設
$C > 0$；$C = 0$ 時三者一樣）。

(c) A、B、C 是正的常數。 foo 的最高次項是 $CN^2$，是 $\Theta(N^2)$（第~\ref{sec:ana-notation}~節
例 3 的多項式只看最高次）。 goo 是 $\frac{C}{2}N^2 + (B + \frac{C}{2})N + A$，也是
$\Theta(N^2)$。 poo 是 $(B + 250C)N + A$，是 $\Theta(N)$。

(d) poo。 它是 $\Theta(N)$，另外兩個是 $\Theta(N^2)$，$N$ 夠大以後一定是 poo 最快。
交叉點也算得出來：poo 的 C 比 foo 少，等於 $250N < N^2$，也就是 $N > 250$；比 goo
少，等於 $250N < N(N+1)/2$，也就是 $N > 499$。 所以 $N \ge 500$ 時 poo 最快，
$N = 499$ 時 poo 和 goo 的 C 都做 $124750$ 次。 (b) 和 (d) 的答案不同，因為漸近的
比較只管 $N$ 夠大的情形，$N = 100$ 時常數 $250$ 還佔上風。 （用小程式實際數過
$N = 100$、$499$、$500$ 時三段程式裡 C 執行的次數，和公式一致。）

\pastsol{px:ana-4}{中山資工所 112 學年度 第 2 題}
題目的版本用相鄰交換：每一輪把 $a[j]$ 和左邊比它大的鄰居交換，一路換到左邊不比它
大為止。 一次交換讓兩個相鄰而且逆序的數變成順序對，其他數對的前後關係不變，所以剛好
消掉一個逆序對。 這和第~\ref{sec:ana-insertion}~節「挪動次數等於逆序對個數」是同一
個論證：\texttt{Swap} 執行的次數等於輸入的逆序對個數。

(a) $9$ 次。 每一輪之後的陣列：
\begin{center}
\begin{tabular}{cccl}
\toprule
$i$ & 這一輪的 $a[i]$ & 交換次數 & 這一輪之後 \\
\midrule
2 & 9 & 0 & 3 9 6 5 8 2 \\
3 & 6 & 1 & 3 6 9 5 8 2 \\
4 & 5 & 2 & 3 5 6 9 8 2 \\
5 & 8 & 1 & 3 5 6 8 9 2 \\
6 & 2 & 5 & 2 3 5 6 8 9 \\
\midrule
合計 & & 9 & \\
\bottomrule
\end{tabular}
\end{center}
用逆序對驗算：$(3,2)$、$(9,6)$、$(9,5)$、$(9,8)$、$(9,2)$、$(6,5)$、$(6,2)$、
$(5,2)$、$(8,2)$，共 $9$ 個。

(b) $1, 2, \ldots, n$（已經排好），$0$ 次。 交換次數等於逆序對個數，不會是負的；等於
$0$ 表示沒有任何逆序對，也就是已經由小到大排好，而 $\{1, \ldots, n\}$ 排好的排列只有
這一個。

(c) $n, n-1, \ldots, 1$（完全相反），$n(n-1)/2$ 次。 $n$ 個位置共有 $\binom{n}{2}$
對，每一對最多貢獻一個逆序對，所以交換次數最多 $n(n-1)/2$；要每一對都是逆序，排列
只能是完全相反的那一個。 直接數也一樣：第 $i$ 輪的 $a[i]$ 比左邊每個數都小，要換
$i - 1$ 次，$\sum_{i=2}^{n} (i - 1) = n(n-1)/2$。 （用程式對 $n \le 7$ 的所有排列檢
查過，最少和最多都只有這兩個排列達到。）

\pastsol{px:ana-5}{政大資科所 110 學年度 第 1 題}
兩句都是錯的。 原題 b 的小寫 $\theta$ 指的就是 $\Theta$。

a.\ 錯。 這句話說的是「每一個排序演算法」，要否定它，舉出一個不是 $O(n^2)$ 的排序
演算法就夠了。 例子：先用遞迴產生 $n$ 個位置的全部 $n!$ 種排列（範例~\ref{ex:avg-shift}
的 \texttt{gen} 就是這樣產生排列），每一種都花 $n - 1$ 次比較檢查照這個順序讀出來是
不是由小到大，最後輸出排好的那一種。 排好的順序一定在這 $n!$ 種裡，所以輸出正確；
不管輸入是什麼，比較次數都是 $n! \cdot (n - 1)$。 而
\[
\frac{n! \cdot (n-1)}{n^2} = \frac{(n-1)! \cdot (n-1)}{n} \ge \frac{(n-1)!}{2} \to \infty
\qquad (n \ge 2),
\]
由定理~\ref{thm:limit}，$n!\,(n-1) \ne O(n^2)$。 這個演算法很笨，但它是一個正確的排
序演算法，所以原句不成立。 如果把原句讀成「排序這個問題可以在 $O(n^2)$ 時間內解決」，
那是對的（插入排序最壞是 $\Theta(n^2)$，第~\ref{sec:ana-cases}~節）；作答時把這個差
別寫出來。

b.\ 錯。 插入排序的最好情況是 $\Theta(n)$。 輸入已經排好時，每一輪第一次比較就停，
共 $n - 1$ 次比較、$0$ 次挪動，第~\ref{sec:ana-line-cost}~節算出執行時間是 $an + b$
的形式，$a > 0$；任何輸入都至少要把 $n$ 個數各看一次，最好情況也不會更少。 而
$n / (n \lg n) = 1 / \lg n \to 0$，由定理~\ref{thm:limit}，$n \ne \Omega(n \lg n)$，
所以最好情況的時間不是 $\Theta(n \log n)$。
